% This file is intended to replace the current chapters/methodology.tex and
% chapters/experimental_setup.tex. In the current article-class template, add the
% top-level section in thesis.tex and input this file there:
% \section{Framework and Experimental Design}
% \label{framework-design}
% % This file is intended to replace the current chapters/methodology.tex and
% chapters/experimental_setup.tex. In the current article-class template, add the
% top-level section in thesis.tex and input this file there:
% \section{Framework and Experimental Design}
% \label{framework-design}
% % This file is intended to replace the current chapters/methodology.tex and
% chapters/experimental_setup.tex. In the current article-class template, add the
% top-level section in thesis.tex and input this file there:
% \section{Framework and Experimental Design}
% \label{framework-design}
% % This file is intended to replace the current chapters/methodology.tex and
% chapters/experimental_setup.tex. In the current article-class template, add the
% top-level section in thesis.tex and input this file there:
% \section{Framework and Experimental Design}
% \label{framework-design}
% \input{chapters/framework_experimental_design}

This chapter defines the analysis framework and its experimental implementation. The
main goal is to make predictive multiplicity observable at the level of individual test
observations and then to test whether this instability is spatially organized in feature
space. The chapter first describes how near-optimal models are collected and how
observation-wise disagreement is measured. It then defines the feature-space graph,
Moran's $I$, LISA hotspot detection, the permutation null model, and the additional
robustness, validation, and attribution analyses used in the empirical study.

\subsection{Overview of the Analysis Pipeline}
\label{sec:framework-overview}

The analysis starts from a supervised binary classification dataset and a finite pool of
trained candidate models. Candidate models are ranked by validation Brier score, and the
main Rashomon approximation is defined as the top-$K$ models in this ranking. For each
test observation, the selected models produce a vector of predicted probabilities. The
primary multiplicity measure is the variance of these predicted probabilities across the
selected models.

The resulting vector of observation-wise predictive variance values is then analyzed as a
spatial signal on a graph. The graph is constructed in the preprocessed feature space of
the test set: observations are nodes, and edges connect nearby observations according to a
symmetrized $k$-nearest-neighbor ($k$NN) relation. Global Moran's $I$ is used to test
whether high-variance observations tend to lie near other high-variance observations.
Local Moran statistics, or Local Indicators of Spatial Association (LISA), are then used
to identify local High--High (HH) regions. These HH regions are interpreted as
multiplicity hotspots: observations with above-average predictive variance surrounded by
neighbors that also have above-average predictive variance.

Finally, the observed spatial structure is compared with a prediction-permutation null
model. This null preserves each model's marginal distribution of predictions but destroys
the alignment between predictions and feature-space location. The empirical analyses then
ask whether the detected hotspot structure is stable across runs, robust to analytical
choices, recoverable in synthetic data, and attributable to model-family or
hyperparameter choices.

\subsection{Data, Model Pool, and Rashomon Approximation}
\label{sec:framework-data-model-pool}

\paragraph{Problem setting and notation.}
Let $(X,y)$ denote a supervised learning dataset with $N$ observations, where
$X \in \mathbb{R}^{N \times d}$ is the feature matrix and
$y \in \{0,1\}^N$ is a binary target variable. Let $\mathcal{F}$ denote a finite pool of
trained probabilistic classifiers, potentially spanning multiple model families and
hyperparameter configurations. Each model $f \in \mathcal{F}$ outputs predicted
probabilities $\hat{p}_{fi} \in [0,1]$ for observations $i=1,\dots,N$.

\paragraph{Datasets.}
Experiments are conducted on three real-world benchmark datasets from different
application domains. All three tasks are binary classification problems and contain both
numerical and categorical input features.

\begin{itemize}
    \item \textbf{COMPAS}: The ProPublica COMPAS recidivism dataset
    \citep{angwin_machine_2016, larson_how_2016}. It contains defendant-level
    information from Broward County, Florida, such as age, prior offence count, sex,
    race, and charge degree. The target indicates whether the individual reoffended
    within two years.

    \item \textbf{German Credit}: The \texttt{credit-g} dataset from OpenML
    \citep{hofmann_statlog_1994}. It contains financial and demographic information
    about credit applicants. The target indicates good versus bad credit risk.

    \item \textbf{Adult}: The Adult income dataset from OpenML
    \citep{becker_adult_1996}. It contains demographic and employment-related
    information. The target indicates whether annual income exceeds \$50{,}000.
\end{itemize}

The datasets differ in size, domain, and feature composition. Additional dataset
characteristics, including split sizes, target prevalence, and transformed feature
dimensionality, are reported in Appendix~\ref{app:dataset-characteristics}.
These differences provide a basic check of whether the proposed spatial multiplicity
analysis behaves similarly across different tabular classification settings.

\paragraph{Data splitting.}
\label{sec:exp-splits}
Each dataset is split into training (60\%), validation (20\%), and test (20\%) sets using
stratified sampling. The training set is used to fit all candidate models. The validation
set is used to rank candidates by Brier score and to construct the finite top-$K$ model
set. The test set is held out until the final analysis step and is used only to compute
predictive variance and spatial statistics.

The full experiment is repeated with 10 random seeds. Two splitting modes are used.
For the main analyses, each seed creates a new stratified train--validation--test split;
this reflects variability in the full pipeline, including model fitting, validation-based
selection, test-set prediction, and spatial analysis. For the cross-seed stability analysis,
the test set is drawn once and kept fixed across seeds while the training and validation
parts are re-split. This fixed-test-set protocol makes pointwise comparisons such as HH
selection frequency and Jaccard overlap meaningful, because the same test observations are
compared across runs.

\paragraph{Model families, training, and preprocessing.}
All models are implemented in Python using \texttt{scikit-learn}
\citep{pedregosa2011scikit}. The candidate pool contains five model families: logistic
regression with elastic-net regularization, $k$-nearest neighbors (kNN), random forest,
gradient boosting, and multilayer perceptron (MLP). These families cover a linear model,
an instance-based method, two tree-based ensemble methods, and a neural-network model.

For each model family and run, 50 hyperparameter configurations are sampled from
predefined finite search spaces. The full search spaces are reported in
Appendix~\ref{app:hp-search-space}. Configurations are sampled using deterministic seeds
derived from the run. If duplicate configurations occur for families with smaller discrete
search spaces, configurations are resampled until 50 unique candidates are obtained. No
cross-validation is used inside a run; instead, the full pipeline is repeated across 10 outer
train--validation--test splits.

Preprocessing is applied uniformly across all model families via a column transformer
\label{sec:exp-preproc}
fitted on the training set only. Numeric features are imputed with the training-set median
and standardized to zero mean and unit variance. Categorical features are imputed with
the training-set mode and one-hot encoded; unknown categories at validation or test time
are mapped to a zero vector. Columns not identified as numeric or categorical are
dropped. The fitted transformer is applied to validation and test data without refitting,
thereby avoiding data leakage. The same transformed feature space is used for model
training and for constructing the feature-space graph.

\paragraph{Rank-based Rashomon approximation.}
\label{sec:meth-rashomon}
\label{sec:exp-rashomon}
The classical Rashomon set is usually defined through a loss tolerance around the best
model \citep{breiman_statistical_2001, marx_predictive_2020,
semenova_existence_2022}. In this thesis, we use a finite, rank-based approximation
instead. The empirical model set is therefore not defined by a fixed absolute or relative
performance tolerance, but by selecting a fixed number of near-optimal models from the
trained candidate pool.

Each model $f \in \mathcal{F}$ is evaluated on the validation set using a loss $L(f)$. In
this thesis, $L$ is the Brier score, because the subsequent multiplicity analysis is based
on predicted probabilities rather than only on hard class labels. Let
$f_{(1)},\ldots,f_{(|\mathcal{F}|)}$ denote the models sorted by increasing validation loss,
so that
\[
L(f_{(1)}) \leq L(f_{(2)}) \leq \cdots \leq L(f_{(|\mathcal{F}|)}).
\]
The main empirical Rashomon approximation is
\[
\mathcal{R}_K = \{f_{(1)},\ldots,f_{(K)}\}.
\]
Unless otherwise stated, $K=25$ and $\mathcal{R}_K$ refers to the global top-$K$ set
selected from the full candidate pool across all families.

The reason for using a fixed $K$ rather than a fixed $\varepsilon$-based performance
threshold is comparability. A fixed tolerance can lead to very different numbers of
selected models across datasets, splits, and model families. This would make
observation-wise variance estimates, spatial hotspot analyses, and comparisons across
runs harder to interpret. Fixing $K$ ensures that each run is based on the same number of
near-optimal models while retaining the central idea of predictive multiplicity: studying
disagreement among models with very similar validation performance.

To audit how close the retained models are in validation performance, we report the
validation Brier-score gap within the selected global top-$K$ set. For each run, this gap
is computed relative to the best retained model,
\[
\Delta L_m = L(f_m) - \min_{f \in \mathcal{R}_K} L(f),
\qquad f_m \in \mathcal{R}_K .
\]
Table~\ref{tab:topk-brier-gap} summarizes the maximum, median, and 95th-percentile
gap across the selected top-$K=25$ models. This diagnostic makes the term
``near-optimal'' explicit for the finite rank-based approximation used in the experiments.

\begin{table}[htbp]
\centering
\caption{Validation Brier-score gaps within the global top-$K=25$ Rashomon
approximation. For each run, gaps are computed relative to the best retained model.
Values are mean $\pm$ standard deviation across runs.}
\label{tab:topk-brier-gap}
\vspace{0.4em}
\input{presentation_assets/tab/topk_brier_gap_summary.tex}
\end{table}

The gaps are small in absolute Brier-score units across all three datasets. German Credit
has the widest selected top-$K$ range, but even there the average maximum gap remains
below $0.01$. This supports interpreting the selected models as a near-optimal finite
candidate set, while making clear that near-optimality is defined operationally by the
rank-based top-$K$ construction.

For analyses that compare model families or study variation within one family, we also
use family-specific top-$K$ sets. If $\mathcal{F}_g \subset \mathcal{F}$ denotes the
candidate models belonging to family $g$, then $\mathcal{R}^{(g)}_K$ denotes the top-$K$
models within that family according to validation Brier score. This construction retains
enough models from each family for the family-wise and within-family attribution
analyses.

\subsection{Measuring Observation-Wise Predictive Multiplicity}
\label{sec:meth-metrics}

For each run, all models in the selected Rashomon approximation are applied to the
held-out test set. This yields a prediction matrix
\[
P \in \mathbb{R}^{K \times N},
\]
where $K$ is the number of selected models and $N$ is the number of test observations.
The entry $\hat{p}_{mi}$ denotes the predicted probability assigned to test observation
$i$ by model $m \in \mathcal{R}_K$.

The primary observation-wise multiplicity metric is predictive variance. Let
$M = |\mathcal{R}_K|$ and let
\[
\bar{p}_i =
\frac{1}{M}
\sum_{m=1}^{M}
\hat{p}_{mi}
\]
denote the mean predicted probability for observation $i$ across the retained models.
The observation-wise predictive variance is computed as the population variance over the
finite selected model set:
\[
v_i =
\frac{1}{M}
\sum_{m=1}^{M}
\left(\hat{p}_{mi} - \bar{p}_i\right)^2,
\quad i=1,\dots,N.
\]
Thus, the denominator is $M$ rather than $M-1$. This convention matches the empirical
pipeline, where the selected top-$K$ model set is treated as the finite Rashomon
approximation under study rather than as a random sample from a larger population.

High values of $v_i$ indicate that near-optimal models assign substantially different
probabilities to observation $i$, while low values indicate agreement. This follows the
score-based view of predictive multiplicity, where disagreement is measured at the level
of predicted scores or probabilities rather than only after thresholding into class labels
\citep{watson-daniels_predictive_2023, hsu_rashomon_nodate}.

An alternative score-based measure is \emph{Rashomon Capacity}
\citep{hsu_rashomon_nodate}. Like predictive variance, it is defined at the level of an
individual observation and captures variation in predicted scores across near-optimal
models. It uses an information-theoretic formulation, whereas predictive variance uses the
empirical spread of predicted probabilities in the finite model set. Predictive variance is
used here because it is directly computable from $\mathcal{R}_K$, easy to interpret as
probability spread, and can be decomposed naturally in the later family- and
hyperparameter-level analyses.

In addition to pointwise predictive variance, the results report aggregate multiplicity
summaries. Mean variance averages $v_i$ across test observations. Ambiguity is the mean
prediction range across Rashomon models,
$\max_m \hat{p}_{mi} - \min_m \hat{p}_{mi}$. The disagreement rate is the fraction of
test observations for which model predictions differ by more than $\delta=0.05$.
Discrepancy is the maximum, over model pairs, of the mean absolute prediction
difference. A supplementary decision-level analysis based on the conflict ratio is reported
in Appendix~\ref{app:conflict-ratio}; it is not part of the main pipeline.

\subsection{Feature-Space Graph and Spatial Hotspot Detection}
\label{sec:meth-graph}

To study whether predictive variance is locally concentrated, the test observations are
represented as a graph in feature space. Each node is one test observation, and edges
connect observations that are similar according to the preprocessed input features.
Throughout this thesis, ``spatial'' therefore refers to neighborhoods in feature space, not
to geographic space.

This section builds on the observation-wise predictive variance defined in
Section~\ref{sec:meth-metrics}, the preprocessing and test-set construction described
in Section~\ref{sec:framework-data-model-pool}, and the feature-space localization motivation introduced
in Section~\ref{sec:spatial-review}.

\paragraph{Graph construction.}
\label{sec:method-spatial}
For each dataset and run, the graph is constructed on the held-out test set. Euclidean
distances are computed between preprocessed test-set feature vectors. For each
observation, its $k$ nearest neighbors are identified, yielding an initial directed $k$NN
graph. Because nearest-neighbor relations are not necessarily symmetric, the graph is
symmetrized: an undirected edge is placed between observations $i$ and $j$ if either
observation is among the $k$ nearest neighbors of the other.

Let
\[
A \in \{0,1\}^{N \times N}
\]
denote the resulting symmetric adjacency matrix, where $A_{ij}=1$ means that
observations $i$ and $j$ are connected. The adjacency matrix is row-standardized to obtain
the spatial weight matrix $W$:
\[
W_{ij} = \frac{A_{ij}}{\sum_{j'=1}^{N} A_{ij'}},
\quad i=1,\dots,N.
\]
After this normalization, each row of $W$ sums to one. Therefore, for an
observation-level quantity $\mathbf{v}=(v_1,\dots,v_N)^\top$, the spatial lag
\[
[W\mathbf{v}]_i = \sum_{j=1}^{N} W_{ij}v_j
\]
is the average value of $\mathbf{v}$ among the neighbors of observation $i$. If $v_i$ is
predictive variance, then $[W\mathbf{v}]_i$ is the average predictive variance in the local
neighborhood of observation $i$.

The baseline graph therefore encodes one specific operational notion of local similarity:
Euclidean distance in the preprocessed feature representation. This choice is natural
because it uses the same transformed feature space as the predictive models, but it is not
unique. In particular, distances in one-hot encoded tabular data can depend on the chosen
preprocessing and scaling. For this reason, the robustness analysis also compares the
baseline graph with PCA-reduced and cosine-distance alternatives
(Section~\ref{sec:results-alt-graph}).

\label{sec:exp-knn-graph}
In the experiments, the reported $k$ values refer to the initial directed nearest-neighbor
query. After symmetrization, individual nodes may have more than $k$ neighbors. The
default values are $k=30$ for COMPAS and German Credit and $k=60$ for Adult, where
the larger value helps avoid disconnected graph components.

\paragraph{Global Moran's $I$.}
\label{sec:meth-moran}
Global Moran's $I$ is a classical measure of spatial autocorrelation \citep{moran1950notes}.
It measures whether similar values tend to occur near each other in a given neighborhood
graph. Here it tests whether high-variance observations tend to be close to other
high-variance observations, and whether low-variance observations tend to be close to
other low-variance observations.

Let $\mathbf{v}=(v_1,\dots,v_N)^\top$ denote the predictive variance values on the test set,
and let $\bar{v}$ be their mean. Because $W$ is row-standardized, Moran's $I$ can be
written as
\[
I =
\frac{\sum_{i=1}^N \sum_{j=1}^N W_{ij}(v_i-\bar{v})(v_j-\bar{v})}
{\sum_{i=1}^N (v_i-\bar{v})^2}.
\]
If observations with above-average variance tend to have neighbors that are also above
average, the numerator becomes positive and Moran's $I$ is positive. Values close to zero
indicate little systematic organization with respect to the graph. Negative values would
indicate that high values tend to be surrounded by low values, or vice versa.

\paragraph{Local Moran statistics and LISA.}
\label{sec:meth-lisa}
Global Moran's $I$ summarizes the overall degree of spatial clustering, but it does not
show where clustering occurs. For this reason, we also use Local Indicators of Spatial
Association (LISA), introduced by \citet{anselin1995local}. LISA decomposes spatial
association into observation-level statistics and identifies observations that contribute to
local clusters. Local Moran's $I$ is the specific LISA statistic used in this thesis; LISA denotes the broader class of local spatial association measures.

In the PySAL implementation (Python Spatial Analysis Library; \citealp{pysal2007}),
the input variable is standardized before local Moran statistics are computed. Let
\[
z_i = \frac{v_i-\bar{v}}{s_v}
\]
denote the standardized predictive variance of observation $i$, where $s_v$ is the standard
deviation of predictive variance across test observations. The spatial lag of the
standardized values is
\[
[W\mathbf{z}]_i = \sum_{j=1}^{N} W_{ij}z_j.
\]
The local Moran statistic used by PySAL is proportional to
\[
z_i[W\mathbf{z}]_i,
\]
up to a constant scaling factor. It is large and positive when observation $i$ and its
neighbors are on the same side of the mean, for example when both have above-average
predictive variance. This connection to the global statistic is the defining LISA property:
up to normalization, the sum of the local Moran statistics is proportional to global
Moran's $I$.

For each observation, PySAL computes a local permutation reference distribution for the
observed Local Moran statistic. This local procedure is conditional on the focal
observation: the value $z_i$ of the focal observation is held fixed, while the values used
to form the local spatial lag around that observation are randomized under the null of
local spatial randomness. The observed local Moran statistic is then compared with this
conditional reference distribution. This yields one local pseudo-$p$-value for each
observation.

Because many local tests are performed at once, some observations may appear significant
by chance alone. We therefore apply the Benjamini--Hochberg procedure
\citep{benjamini_hochberg_1995} to control the false discovery rate before assigning final
LISA labels. Observations that remain significant after this
correction are assigned to one of four local spatial categories:

\begin{itemize}
  \item \textbf{High--High (HH):} $z_i>0$ and $[W\mathbf{z}]_i>0$. The observation has
  above-average predictive variance and is surrounded by neighbors with above-average
  predictive variance. These observations are interpreted as multiplicity hotspots.

  \item \textbf{Low--Low (LL):} $z_i<0$ and $[W\mathbf{z}]_i<0$. The observation has
  below-average predictive variance and is surrounded by neighbors with below-average
  predictive variance. These observations mark locally stable regions.

  \item \textbf{High--Low (HL):} $z_i>0$ and $[W\mathbf{z}]_i<0$. The observation has
  above-average predictive variance but its neighbors have below-average predictive
  variance. This corresponds to a high-variance spatial outlier.

  \item \textbf{Low--High (LH):} $z_i<0$ and $[W\mathbf{z}]_i>0$. The observation has
  below-average predictive variance but its neighbors have above-average predictive
  variance. This corresponds to a low-variance observation inside a high-variance
  neighborhood.
\end{itemize}

Observations that are not significant after correction are labeled non-significant. An
\emph{HH region} refers to the set of High--High observations. An \emph{HH hotspot}
refers to a local region of elevated predictive variance.

\paragraph{Connected HH components.}
The LISA classification identifies individual HH observations. To obtain larger coherent
regions, the symmetrized $k$NN graph is restricted to HH observations only and connected
components are extracted. A connected component is a set of HH observations connected
by paths in the graph. Each component is interpreted as a contiguous high-multiplicity
region in feature space. Very small components may be discarded to avoid treating
isolated fragments as substantive hotspot regions.

For the standard spatial-statistics computations, PySAL's value-permutation procedures
are applied to the already-computed predictive-variance field on the fixed feature-space
graph. Global Moran's $I$ is assessed with global value permutations: the predictive
variance values are randomly reassigned over the graph, producing a reference
distribution for one dataset-level statistic. Local Moran's $I$ is assessed with conditional
local permutations: for each focal observation, the focal value is held fixed while the
surrounding values used for the local comparison are randomized. Thus, the two PySAL
procedures differ in scope. The global test assesses overall spatial autocorrelation,
whereas the local test provides observation-level pseudo-$p$-values used for LISA cluster
labels. These PySAL value-permutation procedures are distinct from the
prediction-matrix permutation null introduced in Section~\ref{sec:meth-null}.

\subsection{Permutation Null Model}
\label{sec:meth-null}
\label{sec:exp-null}

In addition to the standard PySAL value-permutation tests used for Moran's $I$ and LISA,
we use a separate prediction-matrix permutation null. This second null is designed to test
whether the spatial clustering of predictive variance depends on the alignment between
model predictions and feature-space locations.

The null operates on the prediction matrix $P \in \mathbb{R}^{K \times N}$. For each model
$m$ independently, the vector $(\hat{p}_{m1},\dots,\hat{p}_{mN})$ is randomly permuted
across observations. This within-model permutation preserves each model's marginal
distribution of predictions, but destroys systematic alignment between predictions and
the feature-space graph. Predictive variance is then recomputed from the permuted
prediction matrix, and global Moran's $I$ is recalculated. Any spatial clustering in the recomputed variance field is therefore interpreted relative to
a null in which predictions retain their model-specific marginal distributions but are no
longer attached to their original feature-space locations.

For each of $R$ permutations, observation-wise predictive variance and global Moran's
$I$ are recomputed. The empirical one-sided $p$-value is
\[
p = \frac{1 + \#\{\text{null} \geq \text{observed}\}}{R+1}.
\]
In the experiments, $R=100$ null permutations are used per run. A run is treated as
significant when the empirical $p$-value is below 0.05. Because $R=100$, the smallest
attainable empirical $p$-value is $1/(R+1)=1/101\approx 0.0099$. Reported values of
``$p<0.01$'' should therefore be interpreted at this resolution limit.

The prediction-matrix null has a different role from the PySAL value-permutation tests
in Section~\ref{sec:meth-graph}. The PySAL procedures treat the predictive-variance
vector as fixed and test whether those already-computed values are spatially organized
on the fixed graph. The prediction-matrix null instead intervenes one step earlier: it
destroys the alignment between each model's predictions and the feature-space locations,
then recomputes the predictive-variance field. It therefore has power against alternatives
in which model disagreement is systematically tied to regions of feature space. Because
this null deliberately removes the feature-space alignment of the predictions, the
recomputed variance field is expected to be close to spatially unstructured. This also
explains why datasets with strong observed structure can reach the minimum attainable
empirical $p$-value of $1/(R+1)$.

\subsection{Robustness and Validation Design}
\label{sec:framework-robustness-validation}

Several additional analyses test whether the observed spatial structure is robust,
methodologically stable, and recoverable under controlled conditions.

\paragraph{Sensitivity to Rashomon set size.}
The default global Rashomon approximation uses $K=25$. Sensitivity to this choice is examined by varying the Rashomon set size over
\[
K \in \{5, 10, 15, 20, 25, 30, 35, 40, 45, 50\}
\]
using global top-$K$ selection. For each $K$,
predictive variance, Moran's $I$, and HH counts are recomputed.

\paragraph{Sensitivity to neighborhood size.}
The default graph uses dataset-specific $k$ values. To assess whether conclusions depend
on this choice, the neighborhood size is varied and the variance-based spatial metrics are
recomputed. Mean predictive variance itself is independent of $k$, but Moran's $I$ and HH
counts depend on the graph definition.

\paragraph{Calibration robustness.}
\label{sec:exp-calibration}
To test whether the spatial hotspot structure is mainly driven by probability
miscalibration, we apply Platt scaling and isotonic regression to the predictions of each
model in the global Rashomon set. For each model, a logistic regression is fitted on
validation-set predicted probabilities and validation labels, and the resulting mapping is
applied to that model's test-set predictions. The predictive-variance and spatial pipeline
is then rerun on both uncalibrated and calibrated test predictions.

We record changes in mean variance, Moran's $I$, HH/LL masks, connected components,
and the Jaccard similarity between hotspot masks before and after calibration. Calibration
is used here as a post-selection robustness check rather than as an independent
performance evaluation. Because the validation set is used both for selecting the top-$K$
model set and for fitting the calibration mappings, calibrated metrics are interpreted
descriptively. Synthetic datasets are excluded from this analysis because they follow a
separate data-generating design.

\paragraph{Alternative graph constructions.}
To verify that the spatial results are not artifacts of Euclidean distance on the full
one-hot encoded feature space, the baseline graph is compared with two alternatives:
(i) a $k$NN graph in PCA-reduced space retaining 5 components, and (ii) a $k$NN graph
using cosine distance on $\ell_2$-normalized features. All other pipeline parameters are
kept fixed, and Moran's $I$ and HH counts are compared across graph definitions.

\paragraph{Synthetic validation.}
Three synthetic data designs are used to validate whether LISA-based HH regions recover
known ground-truth instability regions under controlled conditions. The single-island
design contains two well-separated stable Gaussian blobs and one ambiguous island, with
points drawn uniformly in a disk and weak XOR-like conditional probability
$P(y=1\mid x)=0.5\pm\delta$. This design validates the pipeline in the classical
aleatoric or boundary-ambiguity setting. The three-islands design extends this setup to
several disconnected ambiguous regions and a small number of outliers with known region
labels; it is reported as an additional validation in Appendix~\ref{app:synthetic-fdr}.
Finally, the structural non-boundary design contains a dense ordinary decision-boundary
strip with $P(y=1\mid x)=0.5$ and two compact exception regions away from that
boundary. This design tests whether the pipeline can recover localized multiplicity when
the intended high-variance regions are not generated by ordinary boundary proximity.
The same training pipeline is applied: a 60/20/20 split, the same five model families,
and top-$K$ selection by validation Brier score. Recovery of the ground-truth regions by
LISA HH classification is evaluated using precision, recall, and Jaccard overlap. The
prediction-permutation null is used to verify that observed Moran's $I$ is significant.

\subsection{Attribution and Interpretation Analyses}
\label{sec:meth-attribution}

The attribution analyses ask which modeling choices are associated with predictive
multiplicity and whether hotspot regions can be described in terms of original features. All
such analyses are descriptive rather than causal: they quantify how disagreement varies
with model family and hyperparameter choices within the sampled candidate pool.

\paragraph{Global and per-family Rashomon sets.}
The global construction selects the top-$K=25$ models from the full candidate pool,
combining all model families. The per-family construction selects the top-$K=25$ models
separately within each family. Comparing these constructions shows whether pooling
across families hides family-specific multiplicity patterns.

\paragraph{Family importance.}
Predictive variance is decomposed into a between-family component and a within-family
component. Let $\mathcal{G}$ denote the set of model families and
$\mathcal{R}^{(g)}_K$ the top-$K$ model set for family $g\in\mathcal{G}$. For observation
$i$, let $\hat{p}_{gmi}$ be the predicted probability from model
$m\in\mathcal{R}^{(g)}_K$. The mean prediction of family $g$ is
\[
\bar{p}_{gi} = \frac{1}{K}\sum_{m\in\mathcal{R}^{(g)}_K}\hat{p}_{gmi},
\]
and the overall mean prediction across family-specific model sets is
\[
\bar{p}_{i} = \frac{1}{|\mathcal{G}|}\sum_{g\in\mathcal{G}}\bar{p}_{gi}.
\]
For a set of observations $S$, such as all test observations, HH observations, or non-HH
observations, the between-family component is
\[
\mathrm{Var}_{\mathrm{between}}(S)
=
\frac{1}{|S|}\sum_{i\in S}
\frac{1}{|\mathcal{G}|}\sum_{g\in\mathcal{G}}
(\bar{p}_{gi}-\bar{p}_{i})^2.
\]
The total predictive variance across all family-specific model sets is
\[
\mathrm{Var}_{\mathrm{total}}(S)
=
\frac{1}{|S|}\sum_{i\in S}
\frac{1}{|\mathcal{G}|K}\sum_{g\in\mathcal{G}}
\sum_{m\in\mathcal{R}^{(g)}_K}(\hat{p}_{gmi}-\bar{p}_{i})^2.
\]
Family importance is defined as
\[
\mathrm{FI}(S) =
\frac{\mathrm{Var}_{\mathrm{between}}(S)}{\mathrm{Var}_{\mathrm{total}}(S)}.
\]
A value close to one means that most observed predictive variance is attributable to
differences between model-family mean predictions. A lower value means that a larger share remains within families,
for example because of hyperparameter variation. Because each family contributes the
same number of models, this decomposition gives equal weight to each model family.

This construction is intentionally balanced across families and is therefore not identical
to the global top-$K$ Rashomon set used to compute the main predictive-variance field
and the LISA HH labels. In the hotspot-specific family-importance analysis, HH membership
is first defined from the global top-$K$ variance field, and the family decomposition is then
evaluated on those same observations using the balanced per-family model sets. The
resulting values should therefore be read as a descriptive comparison of regions, not as a
within-set decomposition of the exact variance field that produced the HH labels.

Because the balanced construction gives each family equal weight, it may emphasize
between-family variation more strongly than an unbalanced empirical Rashomon set in
which some families contribute many near-optimal models and others contribute few.
This is useful for comparing families, but it should not be interpreted as the model-family
share that would arise under every possible Rashomon-set sampling scheme.

\paragraph{Within-family hyperparameter importance.}
After separating the between-family component, the analysis asks which grouped hyperparameter values are associated with disagreement within a given model family. For family $g$, let $\bar{p}^{(g)}_i$ be
the mean prediction across models in $\mathcal{R}^{(g)}_K$:
\[
\bar{p}^{(g)}_i = \frac{1}{K}\sum_{m\in\mathcal{R}^{(g)}_K}\hat{p}^{(g)}_{mi}.
\]
For each model $m\in\mathcal{R}^{(g)}_K$, its model-level disagreement contribution is
\[
V^{(g)}_m =
\frac{1}{N_{\mathrm{test}}}\sum_{i=1}^{N_{\mathrm{test}}}
\left(\hat{p}^{(g)}_{mi}-\bar{p}^{(g)}_i\right)^2.
\]
Thus, $V^{(g)}_m$ measures how strongly model $m$ deviates from the average prediction
of its own family.

For each hyperparameter $h$ within family $g$, models are grouped according to the
observed value or range of that hyperparameter. Categorical hyperparameters and numeric
hyperparameters with at most four distinct values are grouped by exact value. Numeric
hyperparameters with more than four distinct values are discretized into three quantile
bins. If the resulting groups have insufficient support, or fewer than two valid groups
remain, the hyperparameter is excluded for that seed. A one-way variance decomposition
of $V^{(g)}_m$ across the groups of hyperparameter $h$ yields a between-group component
$\mathrm{Var}^{(g)}_{\mathrm{between}}(h)$ and a total component
$\mathrm{Var}^{(g)}_{\mathrm{total}}$. Hyperparameter importance is
\[
\rho^{(g)}(h)
=
\frac{\mathrm{Var}^{(g)}_{\mathrm{between}}(h)}
{\mathrm{Var}^{(g)}_{\mathrm{total}}}.
\]
Results are aggregated over families and random seeds by mean importance, standard
deviation, and the number of valid seeds.

\paragraph{Hotspot-specific hyperparameter shifts.}
To assess whether these descriptive hyperparameter associations differ inside multiplicity hotspots, the same
$V_m$-based importance calculation is repeated on HH observations. For each
hyperparameter, the hotspot-specific importance is compared with its global importance:
\[
\Delta \rho(h) = \rho_{\mathrm{HH}}(h)-\rho_{\mathrm{all}}(h).
\]
Positive values indicate that grouped values of the hyperparameter account for a larger
share of within-family disagreement inside HH regions, while negative values indicate lower importance inside
hotspots. Because HH subsets can be small, especially on German Credit, these estimates
are interpreted descriptively and reported with seed coverage where relevant. An auxiliary
performance-control meta-model diagnostic is reported in Appendix~\ref{app:meta-model-diagnostic}.

\paragraph{Interpretable rule extraction.}
To characterize high-multiplicity regions in terms of original features, shallow decision
trees with maximum depth 3 are used to extract human-readable rules. Trees are fitted on
the original feature space, and positive leaf nodes are converted into conjunctive rules.
The primary target label is HH membership.

Two variants are used. First, global HH rules are fitted to predict membership in the full
HH set of a run. These rules describe broader subgroups enriched for HH observations but
may mix several distinct hotspot regions. Second, component-level rules are fitted for
connected HH components. In this variant, membership in a retained component is the
positive class and all other test observations are the negative class. Component-level rules
therefore aim to describe compact high-multiplicity regions rather than the full HH set at
once.

Rules are evaluated using support, purity, recall, and lift. For global HH rules, purity is
$P(\mathrm{HH}\mid\mathrm{rule\ region})$, and lift is measured relative to the HH base
rate. For component-level rules, purity is
$P(\mathrm{component}\mid\mathrm{rule\ region})$, and lift is measured relative to the
base rate of the corresponding component. The main rule analysis is reported for COMPAS,
where the hotspot structure is clearest and retained components are sufficiently large for
meaningful rule extraction.

\subsection{Chapter Summary}
\label{sec:framework-summary}

\label{sec:meth-assumptions}
The framework relies on four main assumptions. First, the preprocessed feature space must
induce a meaningful notion of local similarity. Second, the finite candidate pool and fixed
top-$K$ selection are treated as a practical approximation of the Rashomon set. Third, the
analysis is descriptive rather than causal: hotspots identify where predictive multiplicity
concentrates, but do not by themselves explain the data-generating mechanisms behind
that disagreement. Fourth, the framework is developed for binary probabilistic
classification and would require adaptation for regression or multiclass settings.

\begin{table}[htbp]
\centering
\small
\caption{Overview of the main analysis goals, methods, and corresponding result sections.}
\label{tab:framework-method-result-map}
\begin{tabular}{p{0.30\textwidth}p{0.34\textwidth}p{0.25\textwidth}}
\toprule
\textbf{Analysis goal} & \textbf{Main method} & \textbf{Result section} \\
\midrule
Measure where near-optimal models disagree & Predictive variance across $\mathcal{R}_K$ & Section~\ref{sec:results-metrics} \\
Test whether disagreement is spatially clustered & Feature-space $k$NN graph, Moran's $I$, prediction-permutation null & Sections~\ref{sec:results-metrics}--\ref{sec:results-null} \\
Identify local high-multiplicity regions & LISA HH labels and connected HH components & Section~\ref{sec:results-spatial} \\
Assess stability and robustness & Fixed-test-set stability, $K$ and $k$ sensitivity, calibration, alternative graphs & Sections~\ref{sec:results-spatial} and~\ref{sec:results-robustness} \\
Summarize associations with modeling choices & Global vs per-family sets, family importance, within-family $V_m$ decomposition & Sections~\ref{sec:results-global-vs-family} and~\ref{sec:results-hp} \\
Validate and interpret hotspots & Synthetic ground-truth designs and shallow-tree rule extraction & Sections~\ref{sec:results-synthetic}--\ref{sec:results-rules} \\
\bottomrule
\end{tabular}
\end{table}


This chapter defines the analysis framework and its experimental implementation. The
main goal is to make predictive multiplicity observable at the level of individual test
observations and then to test whether this instability is spatially organized in feature
space. The chapter first describes how near-optimal models are collected and how
observation-wise disagreement is measured. It then defines the feature-space graph,
Moran's $I$, LISA hotspot detection, the permutation null model, and the additional
robustness, validation, and attribution analyses used in the empirical study.

\subsection{Overview of the Analysis Pipeline}
\label{sec:framework-overview}

The analysis starts from a supervised binary classification dataset and a finite pool of
trained candidate models. Candidate models are ranked by validation Brier score, and the
main Rashomon approximation is defined as the top-$K$ models in this ranking. For each
test observation, the selected models produce a vector of predicted probabilities. The
primary multiplicity measure is the variance of these predicted probabilities across the
selected models.

The resulting vector of observation-wise predictive variance values is then analyzed as a
spatial signal on a graph. The graph is constructed in the preprocessed feature space of
the test set: observations are nodes, and edges connect nearby observations according to a
symmetrized $k$-nearest-neighbor ($k$NN) relation. Global Moran's $I$ is used to test
whether high-variance observations tend to lie near other high-variance observations.
Local Moran statistics, or Local Indicators of Spatial Association (LISA), are then used
to identify local High--High (HH) regions. These HH regions are interpreted as
multiplicity hotspots: observations with above-average predictive variance surrounded by
neighbors that also have above-average predictive variance.

Finally, the observed spatial structure is compared with a prediction-permutation null
model. This null preserves each model's marginal distribution of predictions but destroys
the alignment between predictions and feature-space location. The empirical analyses then
ask whether the detected hotspot structure is stable across runs, robust to analytical
choices, recoverable in synthetic data, and attributable to model-family or
hyperparameter choices.

\subsection{Data, Model Pool, and Rashomon Approximation}
\label{sec:framework-data-model-pool}

\paragraph{Problem setting and notation.}
Let $(X,y)$ denote a supervised learning dataset with $N$ observations, where
$X \in \mathbb{R}^{N \times d}$ is the feature matrix and
$y \in \{0,1\}^N$ is a binary target variable. Let $\mathcal{F}$ denote a finite pool of
trained probabilistic classifiers, potentially spanning multiple model families and
hyperparameter configurations. Each model $f \in \mathcal{F}$ outputs predicted
probabilities $\hat{p}_{fi} \in [0,1]$ for observations $i=1,\dots,N$.

\paragraph{Datasets.}
Experiments are conducted on three real-world benchmark datasets from different
application domains. All three tasks are binary classification problems and contain both
numerical and categorical input features.

\begin{itemize}
    \item \textbf{COMPAS}: The ProPublica COMPAS recidivism dataset
    \citep{angwin_machine_2016, larson_how_2016}. It contains defendant-level
    information from Broward County, Florida, such as age, prior offence count, sex,
    race, and charge degree. The target indicates whether the individual reoffended
    within two years.

    \item \textbf{German Credit}: The \texttt{credit-g} dataset from OpenML
    \citep{hofmann_statlog_1994}. It contains financial and demographic information
    about credit applicants. The target indicates good versus bad credit risk.

    \item \textbf{Adult}: The Adult income dataset from OpenML
    \citep{becker_adult_1996}. It contains demographic and employment-related
    information. The target indicates whether annual income exceeds \$50{,}000.
\end{itemize}

The datasets differ in size, domain, and feature composition. Additional dataset
characteristics, including split sizes, target prevalence, and transformed feature
dimensionality, are reported in Appendix~\ref{app:dataset-characteristics}.
These differences provide a basic check of whether the proposed spatial multiplicity
analysis behaves similarly across different tabular classification settings.

\paragraph{Data splitting.}
\label{sec:exp-splits}
Each dataset is split into training (60\%), validation (20\%), and test (20\%) sets using
stratified sampling. The training set is used to fit all candidate models. The validation
set is used to rank candidates by Brier score and to construct the finite top-$K$ model
set. The test set is held out until the final analysis step and is used only to compute
predictive variance and spatial statistics.

The full experiment is repeated with 10 random seeds. Two splitting modes are used.
For the main analyses, each seed creates a new stratified train--validation--test split;
this reflects variability in the full pipeline, including model fitting, validation-based
selection, test-set prediction, and spatial analysis. For the cross-seed stability analysis,
the test set is drawn once and kept fixed across seeds while the training and validation
parts are re-split. This fixed-test-set protocol makes pointwise comparisons such as HH
selection frequency and Jaccard overlap meaningful, because the same test observations are
compared across runs.

\paragraph{Model families, training, and preprocessing.}
All models are implemented in Python using \texttt{scikit-learn}
\citep{pedregosa2011scikit}. The candidate pool contains five model families: logistic
regression with elastic-net regularization, $k$-nearest neighbors (kNN), random forest,
gradient boosting, and multilayer perceptron (MLP). These families cover a linear model,
an instance-based method, two tree-based ensemble methods, and a neural-network model.

For each model family and run, 50 hyperparameter configurations are sampled from
predefined finite search spaces. The full search spaces are reported in
Appendix~\ref{app:hp-search-space}. Configurations are sampled using deterministic seeds
derived from the run. If duplicate configurations occur for families with smaller discrete
search spaces, configurations are resampled until 50 unique candidates are obtained. No
cross-validation is used inside a run; instead, the full pipeline is repeated across 10 outer
train--validation--test splits.

Preprocessing is applied uniformly across all model families via a column transformer
\label{sec:exp-preproc}
fitted on the training set only. Numeric features are imputed with the training-set median
and standardized to zero mean and unit variance. Categorical features are imputed with
the training-set mode and one-hot encoded; unknown categories at validation or test time
are mapped to a zero vector. Columns not identified as numeric or categorical are
dropped. The fitted transformer is applied to validation and test data without refitting,
thereby avoiding data leakage. The same transformed feature space is used for model
training and for constructing the feature-space graph.

\paragraph{Rank-based Rashomon approximation.}
\label{sec:meth-rashomon}
\label{sec:exp-rashomon}
The classical Rashomon set is usually defined through a loss tolerance around the best
model \citep{breiman_statistical_2001, marx_predictive_2020,
semenova_existence_2022}. In this thesis, we use a finite, rank-based approximation
instead. The empirical model set is therefore not defined by a fixed absolute or relative
performance tolerance, but by selecting a fixed number of near-optimal models from the
trained candidate pool.

Each model $f \in \mathcal{F}$ is evaluated on the validation set using a loss $L(f)$. In
this thesis, $L$ is the Brier score, because the subsequent multiplicity analysis is based
on predicted probabilities rather than only on hard class labels. Let
$f_{(1)},\ldots,f_{(|\mathcal{F}|)}$ denote the models sorted by increasing validation loss,
so that
\[
L(f_{(1)}) \leq L(f_{(2)}) \leq \cdots \leq L(f_{(|\mathcal{F}|)}).
\]
The main empirical Rashomon approximation is
\[
\mathcal{R}_K = \{f_{(1)},\ldots,f_{(K)}\}.
\]
Unless otherwise stated, $K=25$ and $\mathcal{R}_K$ refers to the global top-$K$ set
selected from the full candidate pool across all families.

The reason for using a fixed $K$ rather than a fixed $\varepsilon$-based performance
threshold is comparability. A fixed tolerance can lead to very different numbers of
selected models across datasets, splits, and model families. This would make
observation-wise variance estimates, spatial hotspot analyses, and comparisons across
runs harder to interpret. Fixing $K$ ensures that each run is based on the same number of
near-optimal models while retaining the central idea of predictive multiplicity: studying
disagreement among models with very similar validation performance.

To audit how close the retained models are in validation performance, we report the
validation Brier-score gap within the selected global top-$K$ set. For each run, this gap
is computed relative to the best retained model,
\[
\Delta L_m = L(f_m) - \min_{f \in \mathcal{R}_K} L(f),
\qquad f_m \in \mathcal{R}_K .
\]
Table~\ref{tab:topk-brier-gap} summarizes the maximum, median, and 95th-percentile
gap across the selected top-$K=25$ models. This diagnostic makes the term
``near-optimal'' explicit for the finite rank-based approximation used in the experiments.

\begin{table}[htbp]
\centering
\caption{Validation Brier-score gaps within the global top-$K=25$ Rashomon
approximation. For each run, gaps are computed relative to the best retained model.
Values are mean $\pm$ standard deviation across runs.}
\label{tab:topk-brier-gap}
\vspace{0.4em}
\begin{tabular}{lccc}
\hline
Dataset & Maximum gap & Median gap & 95th-percentile gap \\
\hline
COMPAS & 0.0021 $\pm$ 0.0007 & 0.0013 $\pm$ 0.0005 & 0.0020 $\pm$ 0.0006 \\
German Credit & 0.0076 $\pm$ 0.0025 & 0.0058 $\pm$ 0.0020 & 0.0073 $\pm$ 0.0027 \\
Adult & 0.0039 $\pm$ 0.0006 & 0.0018 $\pm$ 0.0003 & 0.0037 $\pm$ 0.0006 \\
\hline
\end{tabular}
\end{table}

The gaps are small in absolute Brier-score units across all three datasets. German Credit
has the widest selected top-$K$ range, but even there the average maximum gap remains
below $0.01$. This supports interpreting the selected models as a near-optimal finite
candidate set, while making clear that near-optimality is defined operationally by the
rank-based top-$K$ construction.

For analyses that compare model families or study variation within one family, we also
use family-specific top-$K$ sets. If $\mathcal{F}_g \subset \mathcal{F}$ denotes the
candidate models belonging to family $g$, then $\mathcal{R}^{(g)}_K$ denotes the top-$K$
models within that family according to validation Brier score. This construction retains
enough models from each family for the family-wise and within-family attribution
analyses.

\subsection{Measuring Observation-Wise Predictive Multiplicity}
\label{sec:meth-metrics}

For each run, all models in the selected Rashomon approximation are applied to the
held-out test set. This yields a prediction matrix
\[
P \in \mathbb{R}^{K \times N},
\]
where $K$ is the number of selected models and $N$ is the number of test observations.
The entry $\hat{p}_{mi}$ denotes the predicted probability assigned to test observation
$i$ by model $m \in \mathcal{R}_K$.

The primary observation-wise multiplicity metric is predictive variance. Let
$M = |\mathcal{R}_K|$ and let
\[
\bar{p}_i =
\frac{1}{M}
\sum_{m=1}^{M}
\hat{p}_{mi}
\]
denote the mean predicted probability for observation $i$ across the retained models.
The observation-wise predictive variance is computed as the population variance over the
finite selected model set:
\[
v_i =
\frac{1}{M}
\sum_{m=1}^{M}
\left(\hat{p}_{mi} - \bar{p}_i\right)^2,
\quad i=1,\dots,N.
\]
Thus, the denominator is $M$ rather than $M-1$. This convention matches the empirical
pipeline, where the selected top-$K$ model set is treated as the finite Rashomon
approximation under study rather than as a random sample from a larger population.

High values of $v_i$ indicate that near-optimal models assign substantially different
probabilities to observation $i$, while low values indicate agreement. This follows the
score-based view of predictive multiplicity, where disagreement is measured at the level
of predicted scores or probabilities rather than only after thresholding into class labels
\citep{watson-daniels_predictive_2023, hsu_rashomon_nodate}.

An alternative score-based measure is \emph{Rashomon Capacity}
\citep{hsu_rashomon_nodate}. Like predictive variance, it is defined at the level of an
individual observation and captures variation in predicted scores across near-optimal
models. It uses an information-theoretic formulation, whereas predictive variance uses the
empirical spread of predicted probabilities in the finite model set. Predictive variance is
used here because it is directly computable from $\mathcal{R}_K$, easy to interpret as
probability spread, and can be decomposed naturally in the later family- and
hyperparameter-level analyses.

In addition to pointwise predictive variance, the results report aggregate multiplicity
summaries. Mean variance averages $v_i$ across test observations. Ambiguity is the mean
prediction range across Rashomon models,
$\max_m \hat{p}_{mi} - \min_m \hat{p}_{mi}$. The disagreement rate is the fraction of
test observations for which model predictions differ by more than $\delta=0.05$.
Discrepancy is the maximum, over model pairs, of the mean absolute prediction
difference. A supplementary decision-level analysis based on the conflict ratio is reported
in Appendix~\ref{app:conflict-ratio}; it is not part of the main pipeline.

\subsection{Feature-Space Graph and Spatial Hotspot Detection}
\label{sec:meth-graph}

To study whether predictive variance is locally concentrated, the test observations are
represented as a graph in feature space. Each node is one test observation, and edges
connect observations that are similar according to the preprocessed input features.
Throughout this thesis, ``spatial'' therefore refers to neighborhoods in feature space, not
to geographic space.

This section builds on the observation-wise predictive variance defined in
Section~\ref{sec:meth-metrics}, the preprocessing and test-set construction described
in Section~\ref{sec:framework-data-model-pool}, and the feature-space localization motivation introduced
in Section~\ref{sec:spatial-review}.

\paragraph{Graph construction.}
\label{sec:method-spatial}
For each dataset and run, the graph is constructed on the held-out test set. Euclidean
distances are computed between preprocessed test-set feature vectors. For each
observation, its $k$ nearest neighbors are identified, yielding an initial directed $k$NN
graph. Because nearest-neighbor relations are not necessarily symmetric, the graph is
symmetrized: an undirected edge is placed between observations $i$ and $j$ if either
observation is among the $k$ nearest neighbors of the other.

Let
\[
A \in \{0,1\}^{N \times N}
\]
denote the resulting symmetric adjacency matrix, where $A_{ij}=1$ means that
observations $i$ and $j$ are connected. The adjacency matrix is row-standardized to obtain
the spatial weight matrix $W$:
\[
W_{ij} = \frac{A_{ij}}{\sum_{j'=1}^{N} A_{ij'}},
\quad i=1,\dots,N.
\]
After this normalization, each row of $W$ sums to one. Therefore, for an
observation-level quantity $\mathbf{v}=(v_1,\dots,v_N)^\top$, the spatial lag
\[
[W\mathbf{v}]_i = \sum_{j=1}^{N} W_{ij}v_j
\]
is the average value of $\mathbf{v}$ among the neighbors of observation $i$. If $v_i$ is
predictive variance, then $[W\mathbf{v}]_i$ is the average predictive variance in the local
neighborhood of observation $i$.

The baseline graph therefore encodes one specific operational notion of local similarity:
Euclidean distance in the preprocessed feature representation. This choice is natural
because it uses the same transformed feature space as the predictive models, but it is not
unique. In particular, distances in one-hot encoded tabular data can depend on the chosen
preprocessing and scaling. For this reason, the robustness analysis also compares the
baseline graph with PCA-reduced and cosine-distance alternatives
(Section~\ref{sec:results-alt-graph}).

\label{sec:exp-knn-graph}
In the experiments, the reported $k$ values refer to the initial directed nearest-neighbor
query. After symmetrization, individual nodes may have more than $k$ neighbors. The
default values are $k=30$ for COMPAS and German Credit and $k=60$ for Adult, where
the larger value helps avoid disconnected graph components.

\paragraph{Global Moran's $I$.}
\label{sec:meth-moran}
Global Moran's $I$ is a classical measure of spatial autocorrelation \citep{moran1950notes}.
It measures whether similar values tend to occur near each other in a given neighborhood
graph. Here it tests whether high-variance observations tend to be close to other
high-variance observations, and whether low-variance observations tend to be close to
other low-variance observations.

Let $\mathbf{v}=(v_1,\dots,v_N)^\top$ denote the predictive variance values on the test set,
and let $\bar{v}$ be their mean. Because $W$ is row-standardized, Moran's $I$ can be
written as
\[
I =
\frac{\sum_{i=1}^N \sum_{j=1}^N W_{ij}(v_i-\bar{v})(v_j-\bar{v})}
{\sum_{i=1}^N (v_i-\bar{v})^2}.
\]
If observations with above-average variance tend to have neighbors that are also above
average, the numerator becomes positive and Moran's $I$ is positive. Values close to zero
indicate little systematic organization with respect to the graph. Negative values would
indicate that high values tend to be surrounded by low values, or vice versa.

\paragraph{Local Moran statistics and LISA.}
\label{sec:meth-lisa}
Global Moran's $I$ summarizes the overall degree of spatial clustering, but it does not
show where clustering occurs. For this reason, we also use Local Indicators of Spatial
Association (LISA), introduced by \citet{anselin1995local}. LISA decomposes spatial
association into observation-level statistics and identifies observations that contribute to
local clusters. Local Moran's $I$ is the specific LISA statistic used in this thesis; LISA denotes the broader class of local spatial association measures.

In the PySAL implementation (Python Spatial Analysis Library; \citealp{pysal2007}),
the input variable is standardized before local Moran statistics are computed. Let
\[
z_i = \frac{v_i-\bar{v}}{s_v}
\]
denote the standardized predictive variance of observation $i$, where $s_v$ is the standard
deviation of predictive variance across test observations. The spatial lag of the
standardized values is
\[
[W\mathbf{z}]_i = \sum_{j=1}^{N} W_{ij}z_j.
\]
The local Moran statistic used by PySAL is proportional to
\[
z_i[W\mathbf{z}]_i,
\]
up to a constant scaling factor. It is large and positive when observation $i$ and its
neighbors are on the same side of the mean, for example when both have above-average
predictive variance. This connection to the global statistic is the defining LISA property:
up to normalization, the sum of the local Moran statistics is proportional to global
Moran's $I$.

For each observation, PySAL computes a local permutation reference distribution for the
observed Local Moran statistic. This local procedure is conditional on the focal
observation: the value $z_i$ of the focal observation is held fixed, while the values used
to form the local spatial lag around that observation are randomized under the null of
local spatial randomness. The observed local Moran statistic is then compared with this
conditional reference distribution. This yields one local pseudo-$p$-value for each
observation.

Because many local tests are performed at once, some observations may appear significant
by chance alone. We therefore apply the Benjamini--Hochberg procedure
\citep{benjamini_hochberg_1995} to control the false discovery rate before assigning final
LISA labels. Observations that remain significant after this
correction are assigned to one of four local spatial categories:

\begin{itemize}
  \item \textbf{High--High (HH):} $z_i>0$ and $[W\mathbf{z}]_i>0$. The observation has
  above-average predictive variance and is surrounded by neighbors with above-average
  predictive variance. These observations are interpreted as multiplicity hotspots.

  \item \textbf{Low--Low (LL):} $z_i<0$ and $[W\mathbf{z}]_i<0$. The observation has
  below-average predictive variance and is surrounded by neighbors with below-average
  predictive variance. These observations mark locally stable regions.

  \item \textbf{High--Low (HL):} $z_i>0$ and $[W\mathbf{z}]_i<0$. The observation has
  above-average predictive variance but its neighbors have below-average predictive
  variance. This corresponds to a high-variance spatial outlier.

  \item \textbf{Low--High (LH):} $z_i<0$ and $[W\mathbf{z}]_i>0$. The observation has
  below-average predictive variance but its neighbors have above-average predictive
  variance. This corresponds to a low-variance observation inside a high-variance
  neighborhood.
\end{itemize}

Observations that are not significant after correction are labeled non-significant. An
\emph{HH region} refers to the set of High--High observations. An \emph{HH hotspot}
refers to a local region of elevated predictive variance.

\paragraph{Connected HH components.}
The LISA classification identifies individual HH observations. To obtain larger coherent
regions, the symmetrized $k$NN graph is restricted to HH observations only and connected
components are extracted. A connected component is a set of HH observations connected
by paths in the graph. Each component is interpreted as a contiguous high-multiplicity
region in feature space. Very small components may be discarded to avoid treating
isolated fragments as substantive hotspot regions.

For the standard spatial-statistics computations, PySAL's value-permutation procedures
are applied to the already-computed predictive-variance field on the fixed feature-space
graph. Global Moran's $I$ is assessed with global value permutations: the predictive
variance values are randomly reassigned over the graph, producing a reference
distribution for one dataset-level statistic. Local Moran's $I$ is assessed with conditional
local permutations: for each focal observation, the focal value is held fixed while the
surrounding values used for the local comparison are randomized. Thus, the two PySAL
procedures differ in scope. The global test assesses overall spatial autocorrelation,
whereas the local test provides observation-level pseudo-$p$-values used for LISA cluster
labels. These PySAL value-permutation procedures are distinct from the
prediction-matrix permutation null introduced in Section~\ref{sec:meth-null}.

\subsection{Permutation Null Model}
\label{sec:meth-null}
\label{sec:exp-null}

In addition to the standard PySAL value-permutation tests used for Moran's $I$ and LISA,
we use a separate prediction-matrix permutation null. This second null is designed to test
whether the spatial clustering of predictive variance depends on the alignment between
model predictions and feature-space locations.

The null operates on the prediction matrix $P \in \mathbb{R}^{K \times N}$. For each model
$m$ independently, the vector $(\hat{p}_{m1},\dots,\hat{p}_{mN})$ is randomly permuted
across observations. This within-model permutation preserves each model's marginal
distribution of predictions, but destroys systematic alignment between predictions and
the feature-space graph. Predictive variance is then recomputed from the permuted
prediction matrix, and global Moran's $I$ is recalculated. Any spatial clustering in the recomputed variance field is therefore interpreted relative to
a null in which predictions retain their model-specific marginal distributions but are no
longer attached to their original feature-space locations.

For each of $R$ permutations, observation-wise predictive variance and global Moran's
$I$ are recomputed. The empirical one-sided $p$-value is
\[
p = \frac{1 + \#\{\text{null} \geq \text{observed}\}}{R+1}.
\]
In the experiments, $R=100$ null permutations are used per run. A run is treated as
significant when the empirical $p$-value is below 0.05. Because $R=100$, the smallest
attainable empirical $p$-value is $1/(R+1)=1/101\approx 0.0099$. Reported values of
``$p<0.01$'' should therefore be interpreted at this resolution limit.

The prediction-matrix null has a different role from the PySAL value-permutation tests
in Section~\ref{sec:meth-graph}. The PySAL procedures treat the predictive-variance
vector as fixed and test whether those already-computed values are spatially organized
on the fixed graph. The prediction-matrix null instead intervenes one step earlier: it
destroys the alignment between each model's predictions and the feature-space locations,
then recomputes the predictive-variance field. It therefore has power against alternatives
in which model disagreement is systematically tied to regions of feature space. Because
this null deliberately removes the feature-space alignment of the predictions, the
recomputed variance field is expected to be close to spatially unstructured. This also
explains why datasets with strong observed structure can reach the minimum attainable
empirical $p$-value of $1/(R+1)$.

\subsection{Robustness and Validation Design}
\label{sec:framework-robustness-validation}

Several additional analyses test whether the observed spatial structure is robust,
methodologically stable, and recoverable under controlled conditions.

\paragraph{Sensitivity to Rashomon set size.}
The default global Rashomon approximation uses $K=25$. Sensitivity to this choice is examined by varying the Rashomon set size over
\[
K \in \{5, 10, 15, 20, 25, 30, 35, 40, 45, 50\}
\]
using global top-$K$ selection. For each $K$,
predictive variance, Moran's $I$, and HH counts are recomputed.

\paragraph{Sensitivity to neighborhood size.}
The default graph uses dataset-specific $k$ values. To assess whether conclusions depend
on this choice, the neighborhood size is varied and the variance-based spatial metrics are
recomputed. Mean predictive variance itself is independent of $k$, but Moran's $I$ and HH
counts depend on the graph definition.

\paragraph{Calibration robustness.}
\label{sec:exp-calibration}
To test whether the spatial hotspot structure is mainly driven by probability
miscalibration, we apply Platt scaling and isotonic regression to the predictions of each
model in the global Rashomon set. For each model, a logistic regression is fitted on
validation-set predicted probabilities and validation labels, and the resulting mapping is
applied to that model's test-set predictions. The predictive-variance and spatial pipeline
is then rerun on both uncalibrated and calibrated test predictions.

We record changes in mean variance, Moran's $I$, HH/LL masks, connected components,
and the Jaccard similarity between hotspot masks before and after calibration. Calibration
is used here as a post-selection robustness check rather than as an independent
performance evaluation. Because the validation set is used both for selecting the top-$K$
model set and for fitting the calibration mappings, calibrated metrics are interpreted
descriptively. Synthetic datasets are excluded from this analysis because they follow a
separate data-generating design.

\paragraph{Alternative graph constructions.}
To verify that the spatial results are not artifacts of Euclidean distance on the full
one-hot encoded feature space, the baseline graph is compared with two alternatives:
(i) a $k$NN graph in PCA-reduced space retaining 5 components, and (ii) a $k$NN graph
using cosine distance on $\ell_2$-normalized features. All other pipeline parameters are
kept fixed, and Moran's $I$ and HH counts are compared across graph definitions.

\paragraph{Synthetic validation.}
Three synthetic data designs are used to validate whether LISA-based HH regions recover
known ground-truth instability regions under controlled conditions. The single-island
design contains two well-separated stable Gaussian blobs and one ambiguous island, with
points drawn uniformly in a disk and weak XOR-like conditional probability
$P(y=1\mid x)=0.5\pm\delta$. This design validates the pipeline in the classical
aleatoric or boundary-ambiguity setting. The three-islands design extends this setup to
several disconnected ambiguous regions and a small number of outliers with known region
labels; it is reported as an additional validation in Appendix~\ref{app:synthetic-fdr}.
Finally, the structural non-boundary design contains a dense ordinary decision-boundary
strip with $P(y=1\mid x)=0.5$ and two compact exception regions away from that
boundary. This design tests whether the pipeline can recover localized multiplicity when
the intended high-variance regions are not generated by ordinary boundary proximity.
The same training pipeline is applied: a 60/20/20 split, the same five model families,
and top-$K$ selection by validation Brier score. Recovery of the ground-truth regions by
LISA HH classification is evaluated using precision, recall, and Jaccard overlap. The
prediction-permutation null is used to verify that observed Moran's $I$ is significant.

\subsection{Attribution and Interpretation Analyses}
\label{sec:meth-attribution}

The attribution analyses ask which modeling choices are associated with predictive
multiplicity and whether hotspot regions can be described in terms of original features. All
such analyses are descriptive rather than causal: they quantify how disagreement varies
with model family and hyperparameter choices within the sampled candidate pool.

\paragraph{Global and per-family Rashomon sets.}
The global construction selects the top-$K=25$ models from the full candidate pool,
combining all model families. The per-family construction selects the top-$K=25$ models
separately within each family. Comparing these constructions shows whether pooling
across families hides family-specific multiplicity patterns.

\paragraph{Family importance.}
Predictive variance is decomposed into a between-family component and a within-family
component. Let $\mathcal{G}$ denote the set of model families and
$\mathcal{R}^{(g)}_K$ the top-$K$ model set for family $g\in\mathcal{G}$. For observation
$i$, let $\hat{p}_{gmi}$ be the predicted probability from model
$m\in\mathcal{R}^{(g)}_K$. The mean prediction of family $g$ is
\[
\bar{p}_{gi} = \frac{1}{K}\sum_{m\in\mathcal{R}^{(g)}_K}\hat{p}_{gmi},
\]
and the overall mean prediction across family-specific model sets is
\[
\bar{p}_{i} = \frac{1}{|\mathcal{G}|}\sum_{g\in\mathcal{G}}\bar{p}_{gi}.
\]
For a set of observations $S$, such as all test observations, HH observations, or non-HH
observations, the between-family component is
\[
\mathrm{Var}_{\mathrm{between}}(S)
=
\frac{1}{|S|}\sum_{i\in S}
\frac{1}{|\mathcal{G}|}\sum_{g\in\mathcal{G}}
(\bar{p}_{gi}-\bar{p}_{i})^2.
\]
The total predictive variance across all family-specific model sets is
\[
\mathrm{Var}_{\mathrm{total}}(S)
=
\frac{1}{|S|}\sum_{i\in S}
\frac{1}{|\mathcal{G}|K}\sum_{g\in\mathcal{G}}
\sum_{m\in\mathcal{R}^{(g)}_K}(\hat{p}_{gmi}-\bar{p}_{i})^2.
\]
Family importance is defined as
\[
\mathrm{FI}(S) =
\frac{\mathrm{Var}_{\mathrm{between}}(S)}{\mathrm{Var}_{\mathrm{total}}(S)}.
\]
A value close to one means that most observed predictive variance is attributable to
differences between model-family mean predictions. A lower value means that a larger share remains within families,
for example because of hyperparameter variation. Because each family contributes the
same number of models, this decomposition gives equal weight to each model family.

This construction is intentionally balanced across families and is therefore not identical
to the global top-$K$ Rashomon set used to compute the main predictive-variance field
and the LISA HH labels. In the hotspot-specific family-importance analysis, HH membership
is first defined from the global top-$K$ variance field, and the family decomposition is then
evaluated on those same observations using the balanced per-family model sets. The
resulting values should therefore be read as a descriptive comparison of regions, not as a
within-set decomposition of the exact variance field that produced the HH labels.

Because the balanced construction gives each family equal weight, it may emphasize
between-family variation more strongly than an unbalanced empirical Rashomon set in
which some families contribute many near-optimal models and others contribute few.
This is useful for comparing families, but it should not be interpreted as the model-family
share that would arise under every possible Rashomon-set sampling scheme.

\paragraph{Within-family hyperparameter importance.}
After separating the between-family component, the analysis asks which grouped hyperparameter values are associated with disagreement within a given model family. For family $g$, let $\bar{p}^{(g)}_i$ be
the mean prediction across models in $\mathcal{R}^{(g)}_K$:
\[
\bar{p}^{(g)}_i = \frac{1}{K}\sum_{m\in\mathcal{R}^{(g)}_K}\hat{p}^{(g)}_{mi}.
\]
For each model $m\in\mathcal{R}^{(g)}_K$, its model-level disagreement contribution is
\[
V^{(g)}_m =
\frac{1}{N_{\mathrm{test}}}\sum_{i=1}^{N_{\mathrm{test}}}
\left(\hat{p}^{(g)}_{mi}-\bar{p}^{(g)}_i\right)^2.
\]
Thus, $V^{(g)}_m$ measures how strongly model $m$ deviates from the average prediction
of its own family.

For each hyperparameter $h$ within family $g$, models are grouped according to the
observed value or range of that hyperparameter. Categorical hyperparameters and numeric
hyperparameters with at most four distinct values are grouped by exact value. Numeric
hyperparameters with more than four distinct values are discretized into three quantile
bins. If the resulting groups have insufficient support, or fewer than two valid groups
remain, the hyperparameter is excluded for that seed. A one-way variance decomposition
of $V^{(g)}_m$ across the groups of hyperparameter $h$ yields a between-group component
$\mathrm{Var}^{(g)}_{\mathrm{between}}(h)$ and a total component
$\mathrm{Var}^{(g)}_{\mathrm{total}}$. Hyperparameter importance is
\[
\rho^{(g)}(h)
=
\frac{\mathrm{Var}^{(g)}_{\mathrm{between}}(h)}
{\mathrm{Var}^{(g)}_{\mathrm{total}}}.
\]
Results are aggregated over families and random seeds by mean importance, standard
deviation, and the number of valid seeds.

\paragraph{Hotspot-specific hyperparameter shifts.}
To assess whether these descriptive hyperparameter associations differ inside multiplicity hotspots, the same
$V_m$-based importance calculation is repeated on HH observations. For each
hyperparameter, the hotspot-specific importance is compared with its global importance:
\[
\Delta \rho(h) = \rho_{\mathrm{HH}}(h)-\rho_{\mathrm{all}}(h).
\]
Positive values indicate that grouped values of the hyperparameter account for a larger
share of within-family disagreement inside HH regions, while negative values indicate lower importance inside
hotspots. Because HH subsets can be small, especially on German Credit, these estimates
are interpreted descriptively and reported with seed coverage where relevant. An auxiliary
performance-control meta-model diagnostic is reported in Appendix~\ref{app:meta-model-diagnostic}.

\paragraph{Interpretable rule extraction.}
To characterize high-multiplicity regions in terms of original features, shallow decision
trees with maximum depth 3 are used to extract human-readable rules. Trees are fitted on
the original feature space, and positive leaf nodes are converted into conjunctive rules.
The primary target label is HH membership.

Two variants are used. First, global HH rules are fitted to predict membership in the full
HH set of a run. These rules describe broader subgroups enriched for HH observations but
may mix several distinct hotspot regions. Second, component-level rules are fitted for
connected HH components. In this variant, membership in a retained component is the
positive class and all other test observations are the negative class. Component-level rules
therefore aim to describe compact high-multiplicity regions rather than the full HH set at
once.

Rules are evaluated using support, purity, recall, and lift. For global HH rules, purity is
$P(\mathrm{HH}\mid\mathrm{rule\ region})$, and lift is measured relative to the HH base
rate. For component-level rules, purity is
$P(\mathrm{component}\mid\mathrm{rule\ region})$, and lift is measured relative to the
base rate of the corresponding component. The main rule analysis is reported for COMPAS,
where the hotspot structure is clearest and retained components are sufficiently large for
meaningful rule extraction.

\subsection{Chapter Summary}
\label{sec:framework-summary}

\label{sec:meth-assumptions}
The framework relies on four main assumptions. First, the preprocessed feature space must
induce a meaningful notion of local similarity. Second, the finite candidate pool and fixed
top-$K$ selection are treated as a practical approximation of the Rashomon set. Third, the
analysis is descriptive rather than causal: hotspots identify where predictive multiplicity
concentrates, but do not by themselves explain the data-generating mechanisms behind
that disagreement. Fourth, the framework is developed for binary probabilistic
classification and would require adaptation for regression or multiclass settings.

\begin{table}[htbp]
\centering
\small
\caption{Overview of the main analysis goals, methods, and corresponding result sections.}
\label{tab:framework-method-result-map}
\begin{tabular}{p{0.30\textwidth}p{0.34\textwidth}p{0.25\textwidth}}
\toprule
\textbf{Analysis goal} & \textbf{Main method} & \textbf{Result section} \\
\midrule
Measure where near-optimal models disagree & Predictive variance across $\mathcal{R}_K$ & Section~\ref{sec:results-metrics} \\
Test whether disagreement is spatially clustered & Feature-space $k$NN graph, Moran's $I$, prediction-permutation null & Sections~\ref{sec:results-metrics}--\ref{sec:results-null} \\
Identify local high-multiplicity regions & LISA HH labels and connected HH components & Section~\ref{sec:results-spatial} \\
Assess stability and robustness & Fixed-test-set stability, $K$ and $k$ sensitivity, calibration, alternative graphs & Sections~\ref{sec:results-spatial} and~\ref{sec:results-robustness} \\
Summarize associations with modeling choices & Global vs per-family sets, family importance, within-family $V_m$ decomposition & Sections~\ref{sec:results-global-vs-family} and~\ref{sec:results-hp} \\
Validate and interpret hotspots & Synthetic ground-truth designs and shallow-tree rule extraction & Sections~\ref{sec:results-synthetic}--\ref{sec:results-rules} \\
\bottomrule
\end{tabular}
\end{table}


This chapter defines the analysis framework and its experimental implementation. The
main goal is to make predictive multiplicity observable at the level of individual test
observations and then to test whether this instability is spatially organized in feature
space. The chapter first describes how near-optimal models are collected and how
observation-wise disagreement is measured. It then defines the feature-space graph,
Moran's $I$, LISA hotspot detection, the permutation null model, and the additional
robustness, validation, and attribution analyses used in the empirical study.

\subsection{Overview of the Analysis Pipeline}
\label{sec:framework-overview}

The analysis starts from a supervised binary classification dataset and a finite pool of
trained candidate models. Candidate models are ranked by validation Brier score, and the
main Rashomon approximation is defined as the top-$K$ models in this ranking. For each
test observation, the selected models produce a vector of predicted probabilities. The
primary multiplicity measure is the variance of these predicted probabilities across the
selected models.

The resulting vector of observation-wise predictive variance values is then analyzed as a
spatial signal on a graph. The graph is constructed in the preprocessed feature space of
the test set: observations are nodes, and edges connect nearby observations according to a
symmetrized $k$-nearest-neighbor ($k$NN) relation. Global Moran's $I$ is used to test
whether high-variance observations tend to lie near other high-variance observations.
Local Moran statistics, or Local Indicators of Spatial Association (LISA), are then used
to identify local High--High (HH) regions. These HH regions are interpreted as
multiplicity hotspots: observations with above-average predictive variance surrounded by
neighbors that also have above-average predictive variance.

Finally, the observed spatial structure is compared with a prediction-permutation null
model. This null preserves each model's marginal distribution of predictions but destroys
the alignment between predictions and feature-space location. The empirical analyses then
ask whether the detected hotspot structure is stable across runs, robust to analytical
choices, recoverable in synthetic data, and attributable to model-family or
hyperparameter choices.

\subsection{Data, Model Pool, and Rashomon Approximation}
\label{sec:framework-data-model-pool}

\paragraph{Problem setting and notation.}
Let $(X,y)$ denote a supervised learning dataset with $N$ observations, where
$X \in \mathbb{R}^{N \times d}$ is the feature matrix and
$y \in \{0,1\}^N$ is a binary target variable. Let $\mathcal{F}$ denote a finite pool of
trained probabilistic classifiers, potentially spanning multiple model families and
hyperparameter configurations. Each model $f \in \mathcal{F}$ outputs predicted
probabilities $\hat{p}_{fi} \in [0,1]$ for observations $i=1,\dots,N$.

\paragraph{Datasets.}
Experiments are conducted on three real-world benchmark datasets from different
application domains. All three tasks are binary classification problems and contain both
numerical and categorical input features.

\begin{itemize}
    \item \textbf{COMPAS}: The ProPublica COMPAS recidivism dataset
    \citep{angwin_machine_2016, larson_how_2016}. It contains defendant-level
    information from Broward County, Florida, such as age, prior offence count, sex,
    race, and charge degree. The target indicates whether the individual reoffended
    within two years.

    \item \textbf{German Credit}: The \texttt{credit-g} dataset from OpenML
    \citep{hofmann_statlog_1994}. It contains financial and demographic information
    about credit applicants. The target indicates good versus bad credit risk.

    \item \textbf{Adult}: The Adult income dataset from OpenML
    \citep{becker_adult_1996}. It contains demographic and employment-related
    information. The target indicates whether annual income exceeds \$50{,}000.
\end{itemize}

The datasets differ in size, domain, and feature composition. Additional dataset
characteristics, including split sizes, target prevalence, and transformed feature
dimensionality, are reported in Appendix~\ref{app:dataset-characteristics}.
These differences provide a basic check of whether the proposed spatial multiplicity
analysis behaves similarly across different tabular classification settings.

\paragraph{Data splitting.}
\label{sec:exp-splits}
Each dataset is split into training (60\%), validation (20\%), and test (20\%) sets using
stratified sampling. The training set is used to fit all candidate models. The validation
set is used to rank candidates by Brier score and to construct the finite top-$K$ model
set. The test set is held out until the final analysis step and is used only to compute
predictive variance and spatial statistics.

The full experiment is repeated with 10 random seeds. Two splitting modes are used.
For the main analyses, each seed creates a new stratified train--validation--test split;
this reflects variability in the full pipeline, including model fitting, validation-based
selection, test-set prediction, and spatial analysis. For the cross-seed stability analysis,
the test set is drawn once and kept fixed across seeds while the training and validation
parts are re-split. This fixed-test-set protocol makes pointwise comparisons such as HH
selection frequency and Jaccard overlap meaningful, because the same test observations are
compared across runs.

\paragraph{Model families, training, and preprocessing.}
All models are implemented in Python using \texttt{scikit-learn}
\citep{pedregosa2011scikit}. The candidate pool contains five model families: logistic
regression with elastic-net regularization, $k$-nearest neighbors (kNN), random forest,
gradient boosting, and multilayer perceptron (MLP). These families cover a linear model,
an instance-based method, two tree-based ensemble methods, and a neural-network model.

For each model family and run, 50 hyperparameter configurations are sampled from
predefined finite search spaces. The full search spaces are reported in
Appendix~\ref{app:hp-search-space}. Configurations are sampled using deterministic seeds
derived from the run. If duplicate configurations occur for families with smaller discrete
search spaces, configurations are resampled until 50 unique candidates are obtained. No
cross-validation is used inside a run; instead, the full pipeline is repeated across 10 outer
train--validation--test splits.

Preprocessing is applied uniformly across all model families via a column transformer
\label{sec:exp-preproc}
fitted on the training set only. Numeric features are imputed with the training-set median
and standardized to zero mean and unit variance. Categorical features are imputed with
the training-set mode and one-hot encoded; unknown categories at validation or test time
are mapped to a zero vector. Columns not identified as numeric or categorical are
dropped. The fitted transformer is applied to validation and test data without refitting,
thereby avoiding data leakage. The same transformed feature space is used for model
training and for constructing the feature-space graph.

\paragraph{Rank-based Rashomon approximation.}
\label{sec:meth-rashomon}
\label{sec:exp-rashomon}
The classical Rashomon set is usually defined through a loss tolerance around the best
model \citep{breiman_statistical_2001, marx_predictive_2020,
semenova_existence_2022}. In this thesis, we use a finite, rank-based approximation
instead. The empirical model set is therefore not defined by a fixed absolute or relative
performance tolerance, but by selecting a fixed number of near-optimal models from the
trained candidate pool.

Each model $f \in \mathcal{F}$ is evaluated on the validation set using a loss $L(f)$. In
this thesis, $L$ is the Brier score, because the subsequent multiplicity analysis is based
on predicted probabilities rather than only on hard class labels. Let
$f_{(1)},\ldots,f_{(|\mathcal{F}|)}$ denote the models sorted by increasing validation loss,
so that
\[
L(f_{(1)}) \leq L(f_{(2)}) \leq \cdots \leq L(f_{(|\mathcal{F}|)}).
\]
The main empirical Rashomon approximation is
\[
\mathcal{R}_K = \{f_{(1)},\ldots,f_{(K)}\}.
\]
Unless otherwise stated, $K=25$ and $\mathcal{R}_K$ refers to the global top-$K$ set
selected from the full candidate pool across all families.

The reason for using a fixed $K$ rather than a fixed $\varepsilon$-based performance
threshold is comparability. A fixed tolerance can lead to very different numbers of
selected models across datasets, splits, and model families. This would make
observation-wise variance estimates, spatial hotspot analyses, and comparisons across
runs harder to interpret. Fixing $K$ ensures that each run is based on the same number of
near-optimal models while retaining the central idea of predictive multiplicity: studying
disagreement among models with very similar validation performance.

To audit how close the retained models are in validation performance, we report the
validation Brier-score gap within the selected global top-$K$ set. For each run, this gap
is computed relative to the best retained model,
\[
\Delta L_m = L(f_m) - \min_{f \in \mathcal{R}_K} L(f),
\qquad f_m \in \mathcal{R}_K .
\]
Table~\ref{tab:topk-brier-gap} summarizes the maximum, median, and 95th-percentile
gap across the selected top-$K=25$ models. This diagnostic makes the term
``near-optimal'' explicit for the finite rank-based approximation used in the experiments.

\begin{table}[htbp]
\centering
\caption{Validation Brier-score gaps within the global top-$K=25$ Rashomon
approximation. For each run, gaps are computed relative to the best retained model.
Values are mean $\pm$ standard deviation across runs.}
\label{tab:topk-brier-gap}
\vspace{0.4em}
\begin{tabular}{lccc}
\hline
Dataset & Maximum gap & Median gap & 95th-percentile gap \\
\hline
COMPAS & 0.0021 $\pm$ 0.0007 & 0.0013 $\pm$ 0.0005 & 0.0020 $\pm$ 0.0006 \\
German Credit & 0.0076 $\pm$ 0.0025 & 0.0058 $\pm$ 0.0020 & 0.0073 $\pm$ 0.0027 \\
Adult & 0.0039 $\pm$ 0.0006 & 0.0018 $\pm$ 0.0003 & 0.0037 $\pm$ 0.0006 \\
\hline
\end{tabular}
\end{table}

The gaps are small in absolute Brier-score units across all three datasets. German Credit
has the widest selected top-$K$ range, but even there the average maximum gap remains
below $0.01$. This supports interpreting the selected models as a near-optimal finite
candidate set, while making clear that near-optimality is defined operationally by the
rank-based top-$K$ construction.

For analyses that compare model families or study variation within one family, we also
use family-specific top-$K$ sets. If $\mathcal{F}_g \subset \mathcal{F}$ denotes the
candidate models belonging to family $g$, then $\mathcal{R}^{(g)}_K$ denotes the top-$K$
models within that family according to validation Brier score. This construction retains
enough models from each family for the family-wise and within-family attribution
analyses.

\subsection{Measuring Observation-Wise Predictive Multiplicity}
\label{sec:meth-metrics}

For each run, all models in the selected Rashomon approximation are applied to the
held-out test set. This yields a prediction matrix
\[
P \in \mathbb{R}^{K \times N},
\]
where $K$ is the number of selected models and $N$ is the number of test observations.
The entry $\hat{p}_{mi}$ denotes the predicted probability assigned to test observation
$i$ by model $m \in \mathcal{R}_K$.

The primary observation-wise multiplicity metric is predictive variance. Let
$M = |\mathcal{R}_K|$ and let
\[
\bar{p}_i =
\frac{1}{M}
\sum_{m=1}^{M}
\hat{p}_{mi}
\]
denote the mean predicted probability for observation $i$ across the retained models.
The observation-wise predictive variance is computed as the population variance over the
finite selected model set:
\[
v_i =
\frac{1}{M}
\sum_{m=1}^{M}
\left(\hat{p}_{mi} - \bar{p}_i\right)^2,
\quad i=1,\dots,N.
\]
Thus, the denominator is $M$ rather than $M-1$. This convention matches the empirical
pipeline, where the selected top-$K$ model set is treated as the finite Rashomon
approximation under study rather than as a random sample from a larger population.

High values of $v_i$ indicate that near-optimal models assign substantially different
probabilities to observation $i$, while low values indicate agreement. This follows the
score-based view of predictive multiplicity, where disagreement is measured at the level
of predicted scores or probabilities rather than only after thresholding into class labels
\citep{watson-daniels_predictive_2023, hsu_rashomon_nodate}.

An alternative score-based measure is \emph{Rashomon Capacity}
\citep{hsu_rashomon_nodate}. Like predictive variance, it is defined at the level of an
individual observation and captures variation in predicted scores across near-optimal
models. It uses an information-theoretic formulation, whereas predictive variance uses the
empirical spread of predicted probabilities in the finite model set. Predictive variance is
used here because it is directly computable from $\mathcal{R}_K$, easy to interpret as
probability spread, and can be decomposed naturally in the later family- and
hyperparameter-level analyses.

In addition to pointwise predictive variance, the results report aggregate multiplicity
summaries. Mean variance averages $v_i$ across test observations. Ambiguity is the mean
prediction range across Rashomon models,
$\max_m \hat{p}_{mi} - \min_m \hat{p}_{mi}$. The disagreement rate is the fraction of
test observations for which model predictions differ by more than $\delta=0.05$.
Discrepancy is the maximum, over model pairs, of the mean absolute prediction
difference. A supplementary decision-level analysis based on the conflict ratio is reported
in Appendix~\ref{app:conflict-ratio}; it is not part of the main pipeline.

\subsection{Feature-Space Graph and Spatial Hotspot Detection}
\label{sec:meth-graph}

To study whether predictive variance is locally concentrated, the test observations are
represented as a graph in feature space. Each node is one test observation, and edges
connect observations that are similar according to the preprocessed input features.
Throughout this thesis, ``spatial'' therefore refers to neighborhoods in feature space, not
to geographic space.

This section builds on the observation-wise predictive variance defined in
Section~\ref{sec:meth-metrics}, the preprocessing and test-set construction described
in Section~\ref{sec:framework-data-model-pool}, and the feature-space localization motivation introduced
in Section~\ref{sec:spatial-review}.

\paragraph{Graph construction.}
\label{sec:method-spatial}
For each dataset and run, the graph is constructed on the held-out test set. Euclidean
distances are computed between preprocessed test-set feature vectors. For each
observation, its $k$ nearest neighbors are identified, yielding an initial directed $k$NN
graph. Because nearest-neighbor relations are not necessarily symmetric, the graph is
symmetrized: an undirected edge is placed between observations $i$ and $j$ if either
observation is among the $k$ nearest neighbors of the other.

Let
\[
A \in \{0,1\}^{N \times N}
\]
denote the resulting symmetric adjacency matrix, where $A_{ij}=1$ means that
observations $i$ and $j$ are connected. The adjacency matrix is row-standardized to obtain
the spatial weight matrix $W$:
\[
W_{ij} = \frac{A_{ij}}{\sum_{j'=1}^{N} A_{ij'}},
\quad i=1,\dots,N.
\]
After this normalization, each row of $W$ sums to one. Therefore, for an
observation-level quantity $\mathbf{v}=(v_1,\dots,v_N)^\top$, the spatial lag
\[
[W\mathbf{v}]_i = \sum_{j=1}^{N} W_{ij}v_j
\]
is the average value of $\mathbf{v}$ among the neighbors of observation $i$. If $v_i$ is
predictive variance, then $[W\mathbf{v}]_i$ is the average predictive variance in the local
neighborhood of observation $i$.

The baseline graph therefore encodes one specific operational notion of local similarity:
Euclidean distance in the preprocessed feature representation. This choice is natural
because it uses the same transformed feature space as the predictive models, but it is not
unique. In particular, distances in one-hot encoded tabular data can depend on the chosen
preprocessing and scaling. For this reason, the robustness analysis also compares the
baseline graph with PCA-reduced and cosine-distance alternatives
(Section~\ref{sec:results-alt-graph}).

\label{sec:exp-knn-graph}
In the experiments, the reported $k$ values refer to the initial directed nearest-neighbor
query. After symmetrization, individual nodes may have more than $k$ neighbors. The
default values are $k=30$ for COMPAS and German Credit and $k=60$ for Adult, where
the larger value helps avoid disconnected graph components.

\paragraph{Global Moran's $I$.}
\label{sec:meth-moran}
Global Moran's $I$ is a classical measure of spatial autocorrelation \citep{moran1950notes}.
It measures whether similar values tend to occur near each other in a given neighborhood
graph. Here it tests whether high-variance observations tend to be close to other
high-variance observations, and whether low-variance observations tend to be close to
other low-variance observations.

Let $\mathbf{v}=(v_1,\dots,v_N)^\top$ denote the predictive variance values on the test set,
and let $\bar{v}$ be their mean. Because $W$ is row-standardized, Moran's $I$ can be
written as
\[
I =
\frac{\sum_{i=1}^N \sum_{j=1}^N W_{ij}(v_i-\bar{v})(v_j-\bar{v})}
{\sum_{i=1}^N (v_i-\bar{v})^2}.
\]
If observations with above-average variance tend to have neighbors that are also above
average, the numerator becomes positive and Moran's $I$ is positive. Values close to zero
indicate little systematic organization with respect to the graph. Negative values would
indicate that high values tend to be surrounded by low values, or vice versa.

\paragraph{Local Moran statistics and LISA.}
\label{sec:meth-lisa}
Global Moran's $I$ summarizes the overall degree of spatial clustering, but it does not
show where clustering occurs. For this reason, we also use Local Indicators of Spatial
Association (LISA), introduced by \citet{anselin1995local}. LISA decomposes spatial
association into observation-level statistics and identifies observations that contribute to
local clusters. Local Moran's $I$ is the specific LISA statistic used in this thesis; LISA denotes the broader class of local spatial association measures.

In the PySAL implementation (Python Spatial Analysis Library; \citealp{pysal2007}),
the input variable is standardized before local Moran statistics are computed. Let
\[
z_i = \frac{v_i-\bar{v}}{s_v}
\]
denote the standardized predictive variance of observation $i$, where $s_v$ is the standard
deviation of predictive variance across test observations. The spatial lag of the
standardized values is
\[
[W\mathbf{z}]_i = \sum_{j=1}^{N} W_{ij}z_j.
\]
The local Moran statistic used by PySAL is proportional to
\[
z_i[W\mathbf{z}]_i,
\]
up to a constant scaling factor. It is large and positive when observation $i$ and its
neighbors are on the same side of the mean, for example when both have above-average
predictive variance. This connection to the global statistic is the defining LISA property:
up to normalization, the sum of the local Moran statistics is proportional to global
Moran's $I$.

For each observation, PySAL computes a local permutation reference distribution for the
observed Local Moran statistic. This local procedure is conditional on the focal
observation: the value $z_i$ of the focal observation is held fixed, while the values used
to form the local spatial lag around that observation are randomized under the null of
local spatial randomness. The observed local Moran statistic is then compared with this
conditional reference distribution. This yields one local pseudo-$p$-value for each
observation.

Because many local tests are performed at once, some observations may appear significant
by chance alone. We therefore apply the Benjamini--Hochberg procedure
\citep{benjamini_hochberg_1995} to control the false discovery rate before assigning final
LISA labels. Observations that remain significant after this
correction are assigned to one of four local spatial categories:

\begin{itemize}
  \item \textbf{High--High (HH):} $z_i>0$ and $[W\mathbf{z}]_i>0$. The observation has
  above-average predictive variance and is surrounded by neighbors with above-average
  predictive variance. These observations are interpreted as multiplicity hotspots.

  \item \textbf{Low--Low (LL):} $z_i<0$ and $[W\mathbf{z}]_i<0$. The observation has
  below-average predictive variance and is surrounded by neighbors with below-average
  predictive variance. These observations mark locally stable regions.

  \item \textbf{High--Low (HL):} $z_i>0$ and $[W\mathbf{z}]_i<0$. The observation has
  above-average predictive variance but its neighbors have below-average predictive
  variance. This corresponds to a high-variance spatial outlier.

  \item \textbf{Low--High (LH):} $z_i<0$ and $[W\mathbf{z}]_i>0$. The observation has
  below-average predictive variance but its neighbors have above-average predictive
  variance. This corresponds to a low-variance observation inside a high-variance
  neighborhood.
\end{itemize}

Observations that are not significant after correction are labeled non-significant. An
\emph{HH region} refers to the set of High--High observations. An \emph{HH hotspot}
refers to a local region of elevated predictive variance.

\paragraph{Connected HH components.}
The LISA classification identifies individual HH observations. To obtain larger coherent
regions, the symmetrized $k$NN graph is restricted to HH observations only and connected
components are extracted. A connected component is a set of HH observations connected
by paths in the graph. Each component is interpreted as a contiguous high-multiplicity
region in feature space. Very small components may be discarded to avoid treating
isolated fragments as substantive hotspot regions.

For the standard spatial-statistics computations, PySAL's value-permutation procedures
are applied to the already-computed predictive-variance field on the fixed feature-space
graph. Global Moran's $I$ is assessed with global value permutations: the predictive
variance values are randomly reassigned over the graph, producing a reference
distribution for one dataset-level statistic. Local Moran's $I$ is assessed with conditional
local permutations: for each focal observation, the focal value is held fixed while the
surrounding values used for the local comparison are randomized. Thus, the two PySAL
procedures differ in scope. The global test assesses overall spatial autocorrelation,
whereas the local test provides observation-level pseudo-$p$-values used for LISA cluster
labels. These PySAL value-permutation procedures are distinct from the
prediction-matrix permutation null introduced in Section~\ref{sec:meth-null}.

\subsection{Permutation Null Model}
\label{sec:meth-null}
\label{sec:exp-null}

In addition to the standard PySAL value-permutation tests used for Moran's $I$ and LISA,
we use a separate prediction-matrix permutation null. This second null is designed to test
whether the spatial clustering of predictive variance depends on the alignment between
model predictions and feature-space locations.

The null operates on the prediction matrix $P \in \mathbb{R}^{K \times N}$. For each model
$m$ independently, the vector $(\hat{p}_{m1},\dots,\hat{p}_{mN})$ is randomly permuted
across observations. This within-model permutation preserves each model's marginal
distribution of predictions, but destroys systematic alignment between predictions and
the feature-space graph. Predictive variance is then recomputed from the permuted
prediction matrix, and global Moran's $I$ is recalculated. Any spatial clustering in the recomputed variance field is therefore interpreted relative to
a null in which predictions retain their model-specific marginal distributions but are no
longer attached to their original feature-space locations.

For each of $R$ permutations, observation-wise predictive variance and global Moran's
$I$ are recomputed. The empirical one-sided $p$-value is
\[
p = \frac{1 + \#\{\text{null} \geq \text{observed}\}}{R+1}.
\]
In the experiments, $R=100$ null permutations are used per run. A run is treated as
significant when the empirical $p$-value is below 0.05. Because $R=100$, the smallest
attainable empirical $p$-value is $1/(R+1)=1/101\approx 0.0099$. Reported values of
``$p<0.01$'' should therefore be interpreted at this resolution limit.

The prediction-matrix null has a different role from the PySAL value-permutation tests
in Section~\ref{sec:meth-graph}. The PySAL procedures treat the predictive-variance
vector as fixed and test whether those already-computed values are spatially organized
on the fixed graph. The prediction-matrix null instead intervenes one step earlier: it
destroys the alignment between each model's predictions and the feature-space locations,
then recomputes the predictive-variance field. It therefore has power against alternatives
in which model disagreement is systematically tied to regions of feature space. Because
this null deliberately removes the feature-space alignment of the predictions, the
recomputed variance field is expected to be close to spatially unstructured. This also
explains why datasets with strong observed structure can reach the minimum attainable
empirical $p$-value of $1/(R+1)$.

\subsection{Robustness and Validation Design}
\label{sec:framework-robustness-validation}

Several additional analyses test whether the observed spatial structure is robust,
methodologically stable, and recoverable under controlled conditions.

\paragraph{Sensitivity to Rashomon set size.}
The default global Rashomon approximation uses $K=25$. Sensitivity to this choice is examined by varying the Rashomon set size over
\[
K \in \{5, 10, 15, 20, 25, 30, 35, 40, 45, 50\}
\]
using global top-$K$ selection. For each $K$,
predictive variance, Moran's $I$, and HH counts are recomputed.

\paragraph{Sensitivity to neighborhood size.}
The default graph uses dataset-specific $k$ values. To assess whether conclusions depend
on this choice, the neighborhood size is varied and the variance-based spatial metrics are
recomputed. Mean predictive variance itself is independent of $k$, but Moran's $I$ and HH
counts depend on the graph definition.

\paragraph{Calibration robustness.}
\label{sec:exp-calibration}
To test whether the spatial hotspot structure is mainly driven by probability
miscalibration, we apply Platt scaling and isotonic regression to the predictions of each
model in the global Rashomon set. For each model, a logistic regression is fitted on
validation-set predicted probabilities and validation labels, and the resulting mapping is
applied to that model's test-set predictions. The predictive-variance and spatial pipeline
is then rerun on both uncalibrated and calibrated test predictions.

We record changes in mean variance, Moran's $I$, HH/LL masks, connected components,
and the Jaccard similarity between hotspot masks before and after calibration. Calibration
is used here as a post-selection robustness check rather than as an independent
performance evaluation. Because the validation set is used both for selecting the top-$K$
model set and for fitting the calibration mappings, calibrated metrics are interpreted
descriptively. Synthetic datasets are excluded from this analysis because they follow a
separate data-generating design.

\paragraph{Alternative graph constructions.}
To verify that the spatial results are not artifacts of Euclidean distance on the full
one-hot encoded feature space, the baseline graph is compared with two alternatives:
(i) a $k$NN graph in PCA-reduced space retaining 5 components, and (ii) a $k$NN graph
using cosine distance on $\ell_2$-normalized features. All other pipeline parameters are
kept fixed, and Moran's $I$ and HH counts are compared across graph definitions.

\paragraph{Synthetic validation.}
Three synthetic data designs are used to validate whether LISA-based HH regions recover
known ground-truth instability regions under controlled conditions. The single-island
design contains two well-separated stable Gaussian blobs and one ambiguous island, with
points drawn uniformly in a disk and weak XOR-like conditional probability
$P(y=1\mid x)=0.5\pm\delta$. This design validates the pipeline in the classical
aleatoric or boundary-ambiguity setting. The three-islands design extends this setup to
several disconnected ambiguous regions and a small number of outliers with known region
labels; it is reported as an additional validation in Appendix~\ref{app:synthetic-fdr}.
Finally, the structural non-boundary design contains a dense ordinary decision-boundary
strip with $P(y=1\mid x)=0.5$ and two compact exception regions away from that
boundary. This design tests whether the pipeline can recover localized multiplicity when
the intended high-variance regions are not generated by ordinary boundary proximity.
The same training pipeline is applied: a 60/20/20 split, the same five model families,
and top-$K$ selection by validation Brier score. Recovery of the ground-truth regions by
LISA HH classification is evaluated using precision, recall, and Jaccard overlap. The
prediction-permutation null is used to verify that observed Moran's $I$ is significant.

\subsection{Attribution and Interpretation Analyses}
\label{sec:meth-attribution}

The attribution analyses ask which modeling choices are associated with predictive
multiplicity and whether hotspot regions can be described in terms of original features. All
such analyses are descriptive rather than causal: they quantify how disagreement varies
with model family and hyperparameter choices within the sampled candidate pool.

\paragraph{Global and per-family Rashomon sets.}
The global construction selects the top-$K=25$ models from the full candidate pool,
combining all model families. The per-family construction selects the top-$K=25$ models
separately within each family. Comparing these constructions shows whether pooling
across families hides family-specific multiplicity patterns.

\paragraph{Family importance.}
Predictive variance is decomposed into a between-family component and a within-family
component. Let $\mathcal{G}$ denote the set of model families and
$\mathcal{R}^{(g)}_K$ the top-$K$ model set for family $g\in\mathcal{G}$. For observation
$i$, let $\hat{p}_{gmi}$ be the predicted probability from model
$m\in\mathcal{R}^{(g)}_K$. The mean prediction of family $g$ is
\[
\bar{p}_{gi} = \frac{1}{K}\sum_{m\in\mathcal{R}^{(g)}_K}\hat{p}_{gmi},
\]
and the overall mean prediction across family-specific model sets is
\[
\bar{p}_{i} = \frac{1}{|\mathcal{G}|}\sum_{g\in\mathcal{G}}\bar{p}_{gi}.
\]
For a set of observations $S$, such as all test observations, HH observations, or non-HH
observations, the between-family component is
\[
\mathrm{Var}_{\mathrm{between}}(S)
=
\frac{1}{|S|}\sum_{i\in S}
\frac{1}{|\mathcal{G}|}\sum_{g\in\mathcal{G}}
(\bar{p}_{gi}-\bar{p}_{i})^2.
\]
The total predictive variance across all family-specific model sets is
\[
\mathrm{Var}_{\mathrm{total}}(S)
=
\frac{1}{|S|}\sum_{i\in S}
\frac{1}{|\mathcal{G}|K}\sum_{g\in\mathcal{G}}
\sum_{m\in\mathcal{R}^{(g)}_K}(\hat{p}_{gmi}-\bar{p}_{i})^2.
\]
Family importance is defined as
\[
\mathrm{FI}(S) =
\frac{\mathrm{Var}_{\mathrm{between}}(S)}{\mathrm{Var}_{\mathrm{total}}(S)}.
\]
A value close to one means that most observed predictive variance is attributable to
differences between model-family mean predictions. A lower value means that a larger share remains within families,
for example because of hyperparameter variation. Because each family contributes the
same number of models, this decomposition gives equal weight to each model family.

This construction is intentionally balanced across families and is therefore not identical
to the global top-$K$ Rashomon set used to compute the main predictive-variance field
and the LISA HH labels. In the hotspot-specific family-importance analysis, HH membership
is first defined from the global top-$K$ variance field, and the family decomposition is then
evaluated on those same observations using the balanced per-family model sets. The
resulting values should therefore be read as a descriptive comparison of regions, not as a
within-set decomposition of the exact variance field that produced the HH labels.

Because the balanced construction gives each family equal weight, it may emphasize
between-family variation more strongly than an unbalanced empirical Rashomon set in
which some families contribute many near-optimal models and others contribute few.
This is useful for comparing families, but it should not be interpreted as the model-family
share that would arise under every possible Rashomon-set sampling scheme.

\paragraph{Within-family hyperparameter importance.}
After separating the between-family component, the analysis asks which grouped hyperparameter values are associated with disagreement within a given model family. For family $g$, let $\bar{p}^{(g)}_i$ be
the mean prediction across models in $\mathcal{R}^{(g)}_K$:
\[
\bar{p}^{(g)}_i = \frac{1}{K}\sum_{m\in\mathcal{R}^{(g)}_K}\hat{p}^{(g)}_{mi}.
\]
For each model $m\in\mathcal{R}^{(g)}_K$, its model-level disagreement contribution is
\[
V^{(g)}_m =
\frac{1}{N_{\mathrm{test}}}\sum_{i=1}^{N_{\mathrm{test}}}
\left(\hat{p}^{(g)}_{mi}-\bar{p}^{(g)}_i\right)^2.
\]
Thus, $V^{(g)}_m$ measures how strongly model $m$ deviates from the average prediction
of its own family.

For each hyperparameter $h$ within family $g$, models are grouped according to the
observed value or range of that hyperparameter. Categorical hyperparameters and numeric
hyperparameters with at most four distinct values are grouped by exact value. Numeric
hyperparameters with more than four distinct values are discretized into three quantile
bins. If the resulting groups have insufficient support, or fewer than two valid groups
remain, the hyperparameter is excluded for that seed. A one-way variance decomposition
of $V^{(g)}_m$ across the groups of hyperparameter $h$ yields a between-group component
$\mathrm{Var}^{(g)}_{\mathrm{between}}(h)$ and a total component
$\mathrm{Var}^{(g)}_{\mathrm{total}}$. Hyperparameter importance is
\[
\rho^{(g)}(h)
=
\frac{\mathrm{Var}^{(g)}_{\mathrm{between}}(h)}
{\mathrm{Var}^{(g)}_{\mathrm{total}}}.
\]
Results are aggregated over families and random seeds by mean importance, standard
deviation, and the number of valid seeds.

\paragraph{Hotspot-specific hyperparameter shifts.}
To assess whether these descriptive hyperparameter associations differ inside multiplicity hotspots, the same
$V_m$-based importance calculation is repeated on HH observations. For each
hyperparameter, the hotspot-specific importance is compared with its global importance:
\[
\Delta \rho(h) = \rho_{\mathrm{HH}}(h)-\rho_{\mathrm{all}}(h).
\]
Positive values indicate that grouped values of the hyperparameter account for a larger
share of within-family disagreement inside HH regions, while negative values indicate lower importance inside
hotspots. Because HH subsets can be small, especially on German Credit, these estimates
are interpreted descriptively and reported with seed coverage where relevant. An auxiliary
performance-control meta-model diagnostic is reported in Appendix~\ref{app:meta-model-diagnostic}.

\paragraph{Interpretable rule extraction.}
To characterize high-multiplicity regions in terms of original features, shallow decision
trees with maximum depth 3 are used to extract human-readable rules. Trees are fitted on
the original feature space, and positive leaf nodes are converted into conjunctive rules.
The primary target label is HH membership.

Two variants are used. First, global HH rules are fitted to predict membership in the full
HH set of a run. These rules describe broader subgroups enriched for HH observations but
may mix several distinct hotspot regions. Second, component-level rules are fitted for
connected HH components. In this variant, membership in a retained component is the
positive class and all other test observations are the negative class. Component-level rules
therefore aim to describe compact high-multiplicity regions rather than the full HH set at
once.

Rules are evaluated using support, purity, recall, and lift. For global HH rules, purity is
$P(\mathrm{HH}\mid\mathrm{rule\ region})$, and lift is measured relative to the HH base
rate. For component-level rules, purity is
$P(\mathrm{component}\mid\mathrm{rule\ region})$, and lift is measured relative to the
base rate of the corresponding component. The main rule analysis is reported for COMPAS,
where the hotspot structure is clearest and retained components are sufficiently large for
meaningful rule extraction.

\subsection{Chapter Summary}
\label{sec:framework-summary}

\label{sec:meth-assumptions}
The framework relies on four main assumptions. First, the preprocessed feature space must
induce a meaningful notion of local similarity. Second, the finite candidate pool and fixed
top-$K$ selection are treated as a practical approximation of the Rashomon set. Third, the
analysis is descriptive rather than causal: hotspots identify where predictive multiplicity
concentrates, but do not by themselves explain the data-generating mechanisms behind
that disagreement. Fourth, the framework is developed for binary probabilistic
classification and would require adaptation for regression or multiclass settings.

\begin{table}[htbp]
\centering
\small
\caption{Overview of the main analysis goals, methods, and corresponding result sections.}
\label{tab:framework-method-result-map}
\begin{tabular}{p{0.30\textwidth}p{0.34\textwidth}p{0.25\textwidth}}
\toprule
\textbf{Analysis goal} & \textbf{Main method} & \textbf{Result section} \\
\midrule
Measure where near-optimal models disagree & Predictive variance across $\mathcal{R}_K$ & Section~\ref{sec:results-metrics} \\
Test whether disagreement is spatially clustered & Feature-space $k$NN graph, Moran's $I$, prediction-permutation null & Sections~\ref{sec:results-metrics}--\ref{sec:results-null} \\
Identify local high-multiplicity regions & LISA HH labels and connected HH components & Section~\ref{sec:results-spatial} \\
Assess stability and robustness & Fixed-test-set stability, $K$ and $k$ sensitivity, calibration, alternative graphs & Sections~\ref{sec:results-spatial} and~\ref{sec:results-robustness} \\
Summarize associations with modeling choices & Global vs per-family sets, family importance, within-family $V_m$ decomposition & Sections~\ref{sec:results-global-vs-family} and~\ref{sec:results-hp} \\
Validate and interpret hotspots & Synthetic ground-truth designs and shallow-tree rule extraction & Sections~\ref{sec:results-synthetic}--\ref{sec:results-rules} \\
\bottomrule
\end{tabular}
\end{table}


This chapter defines the analysis framework and its experimental implementation. The
main goal is to make predictive multiplicity observable at the level of individual test
observations and then to test whether this instability is spatially organized in feature
space. The chapter first describes how near-optimal models are collected and how
observation-wise disagreement is measured. It then defines the feature-space graph,
Moran's $I$, LISA hotspot detection, the permutation null model, and the additional
robustness, validation, and attribution analyses used in the empirical study.

\subsection{Overview of the Analysis Pipeline}
\label{sec:framework-overview}

The analysis starts from a supervised binary classification dataset and a finite pool of
trained candidate models. Candidate models are ranked by validation Brier score, and the
main Rashomon approximation is defined as the top-$K$ models in this ranking. For each
test observation, the selected models produce a vector of predicted probabilities. The
primary multiplicity measure is the variance of these predicted probabilities across the
selected models.

The resulting vector of observation-wise predictive variance values is then analyzed as a
spatial signal on a graph. The graph is constructed in the preprocessed feature space of
the test set: observations are nodes, and edges connect nearby observations according to a
symmetrized $k$-nearest-neighbor ($k$NN) relation. Global Moran's $I$ is used to test
whether high-variance observations tend to lie near other high-variance observations.
Local Moran statistics, or Local Indicators of Spatial Association (LISA), are then used
to identify local High--High (HH) regions. These HH regions are interpreted as
multiplicity hotspots: observations with above-average predictive variance surrounded by
neighbors that also have above-average predictive variance.

Finally, the observed spatial structure is compared with a prediction-permutation null
model. This null preserves each model's marginal distribution of predictions but destroys
the alignment between predictions and feature-space location. The empirical analyses then
ask whether the detected hotspot structure is stable across runs, robust to analytical
choices, recoverable in synthetic data, and attributable to model-family or
hyperparameter choices.

\subsection{Data, Model Pool, and Rashomon Approximation}
\label{sec:framework-data-model-pool}

\paragraph{Problem setting and notation.}
Let $(X,y)$ denote a supervised learning dataset with $N$ observations, where
$X \in \mathbb{R}^{N \times d}$ is the feature matrix and
$y \in \{0,1\}^N$ is a binary target variable. Let $\mathcal{F}$ denote a finite pool of
trained probabilistic classifiers, potentially spanning multiple model families and
hyperparameter configurations. Each model $f \in \mathcal{F}$ outputs predicted
probabilities $\hat{p}_{fi} \in [0,1]$ for observations $i=1,\dots,N$.

\paragraph{Datasets.}
Experiments are conducted on three real-world benchmark datasets from different
application domains. All three tasks are binary classification problems and contain both
numerical and categorical input features.

\begin{itemize}
    \item \textbf{COMPAS}: The ProPublica COMPAS recidivism dataset
    \citep{angwin_machine_2016, larson_how_2016}. It contains defendant-level
    information from Broward County, Florida, such as age, prior offence count, sex,
    race, and charge degree. The target indicates whether the individual reoffended
    within two years.

    \item \textbf{German Credit}: The \texttt{credit-g} dataset from OpenML
    \citep{hofmann_statlog_1994}. It contains financial and demographic information
    about credit applicants. The target indicates good versus bad credit risk.

    \item \textbf{Adult}: The Adult income dataset from OpenML
    \citep{becker_adult_1996}. It contains demographic and employment-related
    information. The target indicates whether annual income exceeds \$50{,}000.
\end{itemize}

The datasets differ in size, domain, and feature composition. Additional dataset
characteristics, including split sizes, target prevalence, and transformed feature
dimensionality, are reported in Appendix~\ref{app:dataset-characteristics}.
These differences provide a basic check of whether the proposed spatial multiplicity
analysis behaves similarly across different tabular classification settings.

\paragraph{Data splitting.}
\label{sec:exp-splits}
Each dataset is split into training (60\%), validation (20\%), and test (20\%) sets using
stratified sampling. The training set is used to fit all candidate models. The validation
set is used to rank candidates by Brier score and to construct the finite top-$K$ model
set. The test set is held out until the final analysis step and is used only to compute
predictive variance and spatial statistics.

The full experiment is repeated with 10 random seeds. Two splitting modes are used.
For the main analyses, each seed creates a new stratified train--validation--test split;
this reflects variability in the full pipeline, including model fitting, validation-based
selection, test-set prediction, and spatial analysis. For the cross-seed stability analysis,
the test set is drawn once and kept fixed across seeds while the training and validation
parts are re-split. This fixed-test-set protocol makes pointwise comparisons such as HH
selection frequency and Jaccard overlap meaningful, because the same test observations are
compared across runs.

\paragraph{Model families, training, and preprocessing.}
All models are implemented in Python using \texttt{scikit-learn}
\citep{pedregosa2011scikit}. The candidate pool contains five model families: logistic
regression with elastic-net regularization, $k$-nearest neighbors (kNN), random forest,
gradient boosting, and multilayer perceptron (MLP). These families cover a linear model,
an instance-based method, two tree-based ensemble methods, and a neural-network model.

For each model family and run, 50 hyperparameter configurations are sampled from
predefined finite search spaces. The full search spaces are reported in
Appendix~\ref{app:hp-search-space}. Configurations are sampled using deterministic seeds
derived from the run. If duplicate configurations occur for families with smaller discrete
search spaces, configurations are resampled until 50 unique candidates are obtained. No
cross-validation is used inside a run; instead, the full pipeline is repeated across 10 outer
train--validation--test splits.

Preprocessing is applied uniformly across all model families via a column transformer
\label{sec:exp-preproc}
fitted on the training set only. Numeric features are imputed with the training-set median
and standardized to zero mean and unit variance. Categorical features are imputed with
the training-set mode and one-hot encoded; unknown categories at validation or test time
are mapped to a zero vector. Columns not identified as numeric or categorical are
dropped. The fitted transformer is applied to validation and test data without refitting,
thereby avoiding data leakage. The same transformed feature space is used for model
training and for constructing the feature-space graph.

\paragraph{Rank-based Rashomon approximation.}
\label{sec:meth-rashomon}
\label{sec:exp-rashomon}
The classical Rashomon set is usually defined through a loss tolerance around the best
model \citep{breiman_statistical_2001, marx_predictive_2020,
semenova_existence_2022}. In this thesis, we use a finite, rank-based approximation
instead. The empirical model set is therefore not defined by a fixed absolute or relative
performance tolerance, but by selecting a fixed number of near-optimal models from the
trained candidate pool.

Each model $f \in \mathcal{F}$ is evaluated on the validation set using a loss $L(f)$. In
this thesis, $L$ is the Brier score, because the subsequent multiplicity analysis is based
on predicted probabilities rather than only on hard class labels. Let
$f_{(1)},\ldots,f_{(|\mathcal{F}|)}$ denote the models sorted by increasing validation loss,
so that
\[
L(f_{(1)}) \leq L(f_{(2)}) \leq \cdots \leq L(f_{(|\mathcal{F}|)}).
\]
The main empirical Rashomon approximation is
\[
\mathcal{R}_K = \{f_{(1)},\ldots,f_{(K)}\}.
\]
Unless otherwise stated, $K=25$ and $\mathcal{R}_K$ refers to the global top-$K$ set
selected from the full candidate pool across all families.

The reason for using a fixed $K$ rather than a fixed $\varepsilon$-based performance
threshold is comparability. A fixed tolerance can lead to very different numbers of
selected models across datasets, splits, and model families. This would make
observation-wise variance estimates, spatial hotspot analyses, and comparisons across
runs harder to interpret. Fixing $K$ ensures that each run is based on the same number of
near-optimal models while retaining the central idea of predictive multiplicity: studying
disagreement among models with very similar validation performance.

To audit how close the retained models are in validation performance, we report the
validation Brier-score gap within the selected global top-$K$ set. For each run, this gap
is computed relative to the best retained model,
\[
\Delta L_m = L(f_m) - \min_{f \in \mathcal{R}_K} L(f),
\qquad f_m \in \mathcal{R}_K .
\]
Table~\ref{tab:topk-brier-gap} summarizes the maximum, median, and 95th-percentile
gap across the selected top-$K=25$ models. This diagnostic makes the term
``near-optimal'' explicit for the finite rank-based approximation used in the experiments.

\begin{table}[htbp]
\centering
\caption{Validation Brier-score gaps within the global top-$K=25$ Rashomon
approximation. For each run, gaps are computed relative to the best retained model.
Values are mean $\pm$ standard deviation across runs.}
\label{tab:topk-brier-gap}
\vspace{0.4em}
\begin{tabular}{lccc}
\hline
Dataset & Maximum gap & Median gap & 95th-percentile gap \\
\hline
COMPAS & 0.0021 $\pm$ 0.0007 & 0.0013 $\pm$ 0.0005 & 0.0020 $\pm$ 0.0006 \\
German Credit & 0.0076 $\pm$ 0.0025 & 0.0058 $\pm$ 0.0020 & 0.0073 $\pm$ 0.0027 \\
Adult & 0.0039 $\pm$ 0.0006 & 0.0018 $\pm$ 0.0003 & 0.0037 $\pm$ 0.0006 \\
\hline
\end{tabular}
\end{table}

The gaps are small in absolute Brier-score units across all three datasets. German Credit
has the widest selected top-$K$ range, but even there the average maximum gap remains
below $0.01$. This supports interpreting the selected models as a near-optimal finite
candidate set, while making clear that near-optimality is defined operationally by the
rank-based top-$K$ construction.

For analyses that compare model families or study variation within one family, we also
use family-specific top-$K$ sets. If $\mathcal{F}_g \subset \mathcal{F}$ denotes the
candidate models belonging to family $g$, then $\mathcal{R}^{(g)}_K$ denotes the top-$K$
models within that family according to validation Brier score. This construction retains
enough models from each family for the family-wise and within-family attribution
analyses.

\subsection{Measuring Observation-Wise Predictive Multiplicity}
\label{sec:meth-metrics}

For each run, all models in the selected Rashomon approximation are applied to the
held-out test set. This yields a prediction matrix
\[
P \in \mathbb{R}^{K \times N},
\]
where $K$ is the number of selected models and $N$ is the number of test observations.
The entry $\hat{p}_{mi}$ denotes the predicted probability assigned to test observation
$i$ by model $m \in \mathcal{R}_K$.

The primary observation-wise multiplicity metric is predictive variance. Let
$M = |\mathcal{R}_K|$ and let
\[
\bar{p}_i =
\frac{1}{M}
\sum_{m=1}^{M}
\hat{p}_{mi}
\]
denote the mean predicted probability for observation $i$ across the retained models.
The observation-wise predictive variance is computed as the population variance over the
finite selected model set:
\[
v_i =
\frac{1}{M}
\sum_{m=1}^{M}
\left(\hat{p}_{mi} - \bar{p}_i\right)^2,
\quad i=1,\dots,N.
\]
Thus, the denominator is $M$ rather than $M-1$. This convention matches the empirical
pipeline, where the selected top-$K$ model set is treated as the finite Rashomon
approximation under study rather than as a random sample from a larger population.

High values of $v_i$ indicate that near-optimal models assign substantially different
probabilities to observation $i$, while low values indicate agreement. This follows the
score-based view of predictive multiplicity, where disagreement is measured at the level
of predicted scores or probabilities rather than only after thresholding into class labels
\citep{watson-daniels_predictive_2023, hsu_rashomon_nodate}.

An alternative score-based measure is \emph{Rashomon Capacity}
\citep{hsu_rashomon_nodate}. Like predictive variance, it is defined at the level of an
individual observation and captures variation in predicted scores across near-optimal
models. It uses an information-theoretic formulation, whereas predictive variance uses the
empirical spread of predicted probabilities in the finite model set. Predictive variance is
used here because it is directly computable from $\mathcal{R}_K$, easy to interpret as
probability spread, and can be decomposed naturally in the later family- and
hyperparameter-level analyses.

In addition to pointwise predictive variance, the results report aggregate multiplicity
summaries. Mean variance averages $v_i$ across test observations. Ambiguity is the mean
prediction range across Rashomon models,
$\max_m \hat{p}_{mi} - \min_m \hat{p}_{mi}$. The disagreement rate is the fraction of
test observations for which model predictions differ by more than $\delta=0.05$.
Discrepancy is the maximum, over model pairs, of the mean absolute prediction
difference. A supplementary decision-level analysis based on the conflict ratio is reported
in Appendix~\ref{app:conflict-ratio}; it is not part of the main pipeline.

\subsection{Feature-Space Graph and Spatial Hotspot Detection}
\label{sec:meth-graph}

To study whether predictive variance is locally concentrated, the test observations are
represented as a graph in feature space. Each node is one test observation, and edges
connect observations that are similar according to the preprocessed input features.
Throughout this thesis, ``spatial'' therefore refers to neighborhoods in feature space, not
to geographic space.

This section builds on the observation-wise predictive variance defined in
Section~\ref{sec:meth-metrics}, the preprocessing and test-set construction described
in Section~\ref{sec:framework-data-model-pool}, and the feature-space localization motivation introduced
in Section~\ref{sec:spatial-review}.

\paragraph{Graph construction.}
\label{sec:method-spatial}
For each dataset and run, the graph is constructed on the held-out test set. Euclidean
distances are computed between preprocessed test-set feature vectors. For each
observation, its $k$ nearest neighbors are identified, yielding an initial directed $k$NN
graph. Because nearest-neighbor relations are not necessarily symmetric, the graph is
symmetrized: an undirected edge is placed between observations $i$ and $j$ if either
observation is among the $k$ nearest neighbors of the other.

Let
\[
A \in \{0,1\}^{N \times N}
\]
denote the resulting symmetric adjacency matrix, where $A_{ij}=1$ means that
observations $i$ and $j$ are connected. The adjacency matrix is row-standardized to obtain
the spatial weight matrix $W$:
\[
W_{ij} = \frac{A_{ij}}{\sum_{j'=1}^{N} A_{ij'}},
\quad i=1,\dots,N.
\]
After this normalization, each row of $W$ sums to one. Therefore, for an
observation-level quantity $\mathbf{v}=(v_1,\dots,v_N)^\top$, the spatial lag
\[
[W\mathbf{v}]_i = \sum_{j=1}^{N} W_{ij}v_j
\]
is the average value of $\mathbf{v}$ among the neighbors of observation $i$. If $v_i$ is
predictive variance, then $[W\mathbf{v}]_i$ is the average predictive variance in the local
neighborhood of observation $i$.

The baseline graph therefore encodes one specific operational notion of local similarity:
Euclidean distance in the preprocessed feature representation. This choice is natural
because it uses the same transformed feature space as the predictive models, but it is not
unique. In particular, distances in one-hot encoded tabular data can depend on the chosen
preprocessing and scaling. For this reason, the robustness analysis also compares the
baseline graph with PCA-reduced and cosine-distance alternatives
(Section~\ref{sec:results-alt-graph}).

\label{sec:exp-knn-graph}
In the experiments, the reported $k$ values refer to the initial directed nearest-neighbor
query. After symmetrization, individual nodes may have more than $k$ neighbors. The
default values are $k=30$ for COMPAS and German Credit and $k=60$ for Adult, where
the larger value helps avoid disconnected graph components.

\paragraph{Global Moran's $I$.}
\label{sec:meth-moran}
Global Moran's $I$ is a classical measure of spatial autocorrelation \citep{moran1950notes}.
It measures whether similar values tend to occur near each other in a given neighborhood
graph. Here it tests whether high-variance observations tend to be close to other
high-variance observations, and whether low-variance observations tend to be close to
other low-variance observations.

Let $\mathbf{v}=(v_1,\dots,v_N)^\top$ denote the predictive variance values on the test set,
and let $\bar{v}$ be their mean. Because $W$ is row-standardized, Moran's $I$ can be
written as
\[
I =
\frac{\sum_{i=1}^N \sum_{j=1}^N W_{ij}(v_i-\bar{v})(v_j-\bar{v})}
{\sum_{i=1}^N (v_i-\bar{v})^2}.
\]
If observations with above-average variance tend to have neighbors that are also above
average, the numerator becomes positive and Moran's $I$ is positive. Values close to zero
indicate little systematic organization with respect to the graph. Negative values would
indicate that high values tend to be surrounded by low values, or vice versa.

\paragraph{Local Moran statistics and LISA.}
\label{sec:meth-lisa}
Global Moran's $I$ summarizes the overall degree of spatial clustering, but it does not
show where clustering occurs. For this reason, we also use Local Indicators of Spatial
Association (LISA), introduced by \citet{anselin1995local}. LISA decomposes spatial
association into observation-level statistics and identifies observations that contribute to
local clusters. Local Moran's $I$ is the specific LISA statistic used in this thesis; LISA denotes the broader class of local spatial association measures.

In the PySAL implementation (Python Spatial Analysis Library; \citealp{pysal2007}),
the input variable is standardized before local Moran statistics are computed. Let
\[
z_i = \frac{v_i-\bar{v}}{s_v}
\]
denote the standardized predictive variance of observation $i$, where $s_v$ is the standard
deviation of predictive variance across test observations. The spatial lag of the
standardized values is
\[
[W\mathbf{z}]_i = \sum_{j=1}^{N} W_{ij}z_j.
\]
The local Moran statistic used by PySAL is proportional to
\[
z_i[W\mathbf{z}]_i,
\]
up to a constant scaling factor. It is large and positive when observation $i$ and its
neighbors are on the same side of the mean, for example when both have above-average
predictive variance. This connection to the global statistic is the defining LISA property:
up to normalization, the sum of the local Moran statistics is proportional to global
Moran's $I$.

For each observation, PySAL computes a local permutation reference distribution for the
observed Local Moran statistic. This local procedure is conditional on the focal
observation: the value $z_i$ of the focal observation is held fixed, while the values used
to form the local spatial lag around that observation are randomized under the null of
local spatial randomness. The observed local Moran statistic is then compared with this
conditional reference distribution. This yields one local pseudo-$p$-value for each
observation.

Because many local tests are performed at once, some observations may appear significant
by chance alone. We therefore apply the Benjamini--Hochberg procedure
\citep{benjamini_hochberg_1995} to control the false discovery rate before assigning final
LISA labels. Observations that remain significant after this
correction are assigned to one of four local spatial categories:

\begin{itemize}
  \item \textbf{High--High (HH):} $z_i>0$ and $[W\mathbf{z}]_i>0$. The observation has
  above-average predictive variance and is surrounded by neighbors with above-average
  predictive variance. These observations are interpreted as multiplicity hotspots.

  \item \textbf{Low--Low (LL):} $z_i<0$ and $[W\mathbf{z}]_i<0$. The observation has
  below-average predictive variance and is surrounded by neighbors with below-average
  predictive variance. These observations mark locally stable regions.

  \item \textbf{High--Low (HL):} $z_i>0$ and $[W\mathbf{z}]_i<0$. The observation has
  above-average predictive variance but its neighbors have below-average predictive
  variance. This corresponds to a high-variance spatial outlier.

  \item \textbf{Low--High (LH):} $z_i<0$ and $[W\mathbf{z}]_i>0$. The observation has
  below-average predictive variance but its neighbors have above-average predictive
  variance. This corresponds to a low-variance observation inside a high-variance
  neighborhood.
\end{itemize}

Observations that are not significant after correction are labeled non-significant. An
\emph{HH region} refers to the set of High--High observations. An \emph{HH hotspot}
refers to a local region of elevated predictive variance.

\paragraph{Connected HH components.}
The LISA classification identifies individual HH observations. To obtain larger coherent
regions, the symmetrized $k$NN graph is restricted to HH observations only and connected
components are extracted. A connected component is a set of HH observations connected
by paths in the graph. Each component is interpreted as a contiguous high-multiplicity
region in feature space. Very small components may be discarded to avoid treating
isolated fragments as substantive hotspot regions.

For the standard spatial-statistics computations, PySAL's value-permutation procedures
are applied to the already-computed predictive-variance field on the fixed feature-space
graph. Global Moran's $I$ is assessed with global value permutations: the predictive
variance values are randomly reassigned over the graph, producing a reference
distribution for one dataset-level statistic. Local Moran's $I$ is assessed with conditional
local permutations: for each focal observation, the focal value is held fixed while the
surrounding values used for the local comparison are randomized. Thus, the two PySAL
procedures differ in scope. The global test assesses overall spatial autocorrelation,
whereas the local test provides observation-level pseudo-$p$-values used for LISA cluster
labels. These PySAL value-permutation procedures are distinct from the
prediction-matrix permutation null introduced in Section~\ref{sec:meth-null}.

\subsection{Permutation Null Model}
\label{sec:meth-null}
\label{sec:exp-null}

In addition to the standard PySAL value-permutation tests used for Moran's $I$ and LISA,
we use a separate prediction-matrix permutation null. This second null is designed to test
whether the spatial clustering of predictive variance depends on the alignment between
model predictions and feature-space locations.

The null operates on the prediction matrix $P \in \mathbb{R}^{K \times N}$. For each model
$m$ independently, the vector $(\hat{p}_{m1},\dots,\hat{p}_{mN})$ is randomly permuted
across observations. This within-model permutation preserves each model's marginal
distribution of predictions, but destroys systematic alignment between predictions and
the feature-space graph. Predictive variance is then recomputed from the permuted
prediction matrix, and global Moran's $I$ is recalculated. Any spatial clustering in the recomputed variance field is therefore interpreted relative to
a null in which predictions retain their model-specific marginal distributions but are no
longer attached to their original feature-space locations.

For each of $R$ permutations, observation-wise predictive variance and global Moran's
$I$ are recomputed. The empirical one-sided $p$-value is
\[
p = \frac{1 + \#\{\text{null} \geq \text{observed}\}}{R+1}.
\]
In the experiments, $R=100$ null permutations are used per run. A run is treated as
significant when the empirical $p$-value is below 0.05. Because $R=100$, the smallest
attainable empirical $p$-value is $1/(R+1)=1/101\approx 0.0099$. Reported values of
``$p<0.01$'' should therefore be interpreted at this resolution limit.

The prediction-matrix null has a different role from the PySAL value-permutation tests
in Section~\ref{sec:meth-graph}. The PySAL procedures treat the predictive-variance
vector as fixed and test whether those already-computed values are spatially organized
on the fixed graph. The prediction-matrix null instead intervenes one step earlier: it
destroys the alignment between each model's predictions and the feature-space locations,
then recomputes the predictive-variance field. It therefore has power against alternatives
in which model disagreement is systematically tied to regions of feature space. Because
this null deliberately removes the feature-space alignment of the predictions, the
recomputed variance field is expected to be close to spatially unstructured. This also
explains why datasets with strong observed structure can reach the minimum attainable
empirical $p$-value of $1/(R+1)$.

\subsection{Robustness and Validation Design}
\label{sec:framework-robustness-validation}

Several additional analyses test whether the observed spatial structure is robust,
methodologically stable, and recoverable under controlled conditions.

\paragraph{Sensitivity to Rashomon set size.}
The default global Rashomon approximation uses $K=25$. Sensitivity to this choice is examined by varying the Rashomon set size over
\[
K \in \{5, 10, 15, 20, 25, 30, 35, 40, 45, 50\}
\]
using global top-$K$ selection. For each $K$,
predictive variance, Moran's $I$, and HH counts are recomputed.

\paragraph{Sensitivity to neighborhood size.}
The default graph uses dataset-specific $k$ values. To assess whether conclusions depend
on this choice, the neighborhood size is varied and the variance-based spatial metrics are
recomputed. Mean predictive variance itself is independent of $k$, but Moran's $I$ and HH
counts depend on the graph definition.

\paragraph{Calibration robustness.}
\label{sec:exp-calibration}
To test whether the spatial hotspot structure is mainly driven by probability
miscalibration, we apply Platt scaling and isotonic regression to the predictions of each
model in the global Rashomon set. For each model, a logistic regression is fitted on
validation-set predicted probabilities and validation labels, and the resulting mapping is
applied to that model's test-set predictions. The predictive-variance and spatial pipeline
is then rerun on both uncalibrated and calibrated test predictions.

We record changes in mean variance, Moran's $I$, HH/LL masks, connected components,
and the Jaccard similarity between hotspot masks before and after calibration. Calibration
is used here as a post-selection robustness check rather than as an independent
performance evaluation. Because the validation set is used both for selecting the top-$K$
model set and for fitting the calibration mappings, calibrated metrics are interpreted
descriptively. Synthetic datasets are excluded from this analysis because they follow a
separate data-generating design.

\paragraph{Alternative graph constructions.}
To verify that the spatial results are not artifacts of Euclidean distance on the full
one-hot encoded feature space, the baseline graph is compared with two alternatives:
(i) a $k$NN graph in PCA-reduced space retaining 5 components, and (ii) a $k$NN graph
using cosine distance on $\ell_2$-normalized features. All other pipeline parameters are
kept fixed, and Moran's $I$ and HH counts are compared across graph definitions.

\paragraph{Synthetic validation.}
Three synthetic data designs are used to validate whether LISA-based HH regions recover
known ground-truth instability regions under controlled conditions. The single-island
design contains two well-separated stable Gaussian blobs and one ambiguous island, with
points drawn uniformly in a disk and weak XOR-like conditional probability
$P(y=1\mid x)=0.5\pm\delta$. This design validates the pipeline in the classical
aleatoric or boundary-ambiguity setting. The three-islands design extends this setup to
several disconnected ambiguous regions and a small number of outliers with known region
labels; it is reported as an additional validation in Appendix~\ref{app:synthetic-fdr}.
Finally, the structural non-boundary design contains a dense ordinary decision-boundary
strip with $P(y=1\mid x)=0.5$ and two compact exception regions away from that
boundary. This design tests whether the pipeline can recover localized multiplicity when
the intended high-variance regions are not generated by ordinary boundary proximity.
The same training pipeline is applied: a 60/20/20 split, the same five model families,
and top-$K$ selection by validation Brier score. Recovery of the ground-truth regions by
LISA HH classification is evaluated using precision, recall, and Jaccard overlap. The
prediction-permutation null is used to verify that observed Moran's $I$ is significant.

\subsection{Attribution and Interpretation Analyses}
\label{sec:meth-attribution}

The attribution analyses ask which modeling choices are associated with predictive
multiplicity and whether hotspot regions can be described in terms of original features. All
such analyses are descriptive rather than causal: they quantify how disagreement varies
with model family and hyperparameter choices within the sampled candidate pool.

\paragraph{Global and per-family Rashomon sets.}
The global construction selects the top-$K=25$ models from the full candidate pool,
combining all model families. The per-family construction selects the top-$K=25$ models
separately within each family. Comparing these constructions shows whether pooling
across families hides family-specific multiplicity patterns.

\paragraph{Family importance.}
Predictive variance is decomposed into a between-family component and a within-family
component. Let $\mathcal{G}$ denote the set of model families and
$\mathcal{R}^{(g)}_K$ the top-$K$ model set for family $g\in\mathcal{G}$. For observation
$i$, let $\hat{p}_{gmi}$ be the predicted probability from model
$m\in\mathcal{R}^{(g)}_K$. The mean prediction of family $g$ is
\[
\bar{p}_{gi} = \frac{1}{K}\sum_{m\in\mathcal{R}^{(g)}_K}\hat{p}_{gmi},
\]
and the overall mean prediction across family-specific model sets is
\[
\bar{p}_{i} = \frac{1}{|\mathcal{G}|}\sum_{g\in\mathcal{G}}\bar{p}_{gi}.
\]
For a set of observations $S$, such as all test observations, HH observations, or non-HH
observations, the between-family component is
\[
\mathrm{Var}_{\mathrm{between}}(S)
=
\frac{1}{|S|}\sum_{i\in S}
\frac{1}{|\mathcal{G}|}\sum_{g\in\mathcal{G}}
(\bar{p}_{gi}-\bar{p}_{i})^2.
\]
The total predictive variance across all family-specific model sets is
\[
\mathrm{Var}_{\mathrm{total}}(S)
=
\frac{1}{|S|}\sum_{i\in S}
\frac{1}{|\mathcal{G}|K}\sum_{g\in\mathcal{G}}
\sum_{m\in\mathcal{R}^{(g)}_K}(\hat{p}_{gmi}-\bar{p}_{i})^2.
\]
Family importance is defined as
\[
\mathrm{FI}(S) =
\frac{\mathrm{Var}_{\mathrm{between}}(S)}{\mathrm{Var}_{\mathrm{total}}(S)}.
\]
A value close to one means that most observed predictive variance is attributable to
differences between model-family mean predictions. A lower value means that a larger share remains within families,
for example because of hyperparameter variation. Because each family contributes the
same number of models, this decomposition gives equal weight to each model family.

This construction is intentionally balanced across families and is therefore not identical
to the global top-$K$ Rashomon set used to compute the main predictive-variance field
and the LISA HH labels. In the hotspot-specific family-importance analysis, HH membership
is first defined from the global top-$K$ variance field, and the family decomposition is then
evaluated on those same observations using the balanced per-family model sets. The
resulting values should therefore be read as a descriptive comparison of regions, not as a
within-set decomposition of the exact variance field that produced the HH labels.

Because the balanced construction gives each family equal weight, it may emphasize
between-family variation more strongly than an unbalanced empirical Rashomon set in
which some families contribute many near-optimal models and others contribute few.
This is useful for comparing families, but it should not be interpreted as the model-family
share that would arise under every possible Rashomon-set sampling scheme.

\paragraph{Within-family hyperparameter importance.}
After separating the between-family component, the analysis asks which grouped hyperparameter values are associated with disagreement within a given model family. For family $g$, let $\bar{p}^{(g)}_i$ be
the mean prediction across models in $\mathcal{R}^{(g)}_K$:
\[
\bar{p}^{(g)}_i = \frac{1}{K}\sum_{m\in\mathcal{R}^{(g)}_K}\hat{p}^{(g)}_{mi}.
\]
For each model $m\in\mathcal{R}^{(g)}_K$, its model-level disagreement contribution is
\[
V^{(g)}_m =
\frac{1}{N_{\mathrm{test}}}\sum_{i=1}^{N_{\mathrm{test}}}
\left(\hat{p}^{(g)}_{mi}-\bar{p}^{(g)}_i\right)^2.
\]
Thus, $V^{(g)}_m$ measures how strongly model $m$ deviates from the average prediction
of its own family.

For each hyperparameter $h$ within family $g$, models are grouped according to the
observed value or range of that hyperparameter. Categorical hyperparameters and numeric
hyperparameters with at most four distinct values are grouped by exact value. Numeric
hyperparameters with more than four distinct values are discretized into three quantile
bins. If the resulting groups have insufficient support, or fewer than two valid groups
remain, the hyperparameter is excluded for that seed. A one-way variance decomposition
of $V^{(g)}_m$ across the groups of hyperparameter $h$ yields a between-group component
$\mathrm{Var}^{(g)}_{\mathrm{between}}(h)$ and a total component
$\mathrm{Var}^{(g)}_{\mathrm{total}}$. Hyperparameter importance is
\[
\rho^{(g)}(h)
=
\frac{\mathrm{Var}^{(g)}_{\mathrm{between}}(h)}
{\mathrm{Var}^{(g)}_{\mathrm{total}}}.
\]
Results are aggregated over families and random seeds by mean importance, standard
deviation, and the number of valid seeds.

\paragraph{Hotspot-specific hyperparameter shifts.}
To assess whether these descriptive hyperparameter associations differ inside multiplicity hotspots, the same
$V_m$-based importance calculation is repeated on HH observations. For each
hyperparameter, the hotspot-specific importance is compared with its global importance:
\[
\Delta \rho(h) = \rho_{\mathrm{HH}}(h)-\rho_{\mathrm{all}}(h).
\]
Positive values indicate that grouped values of the hyperparameter account for a larger
share of within-family disagreement inside HH regions, while negative values indicate lower importance inside
hotspots. Because HH subsets can be small, especially on German Credit, these estimates
are interpreted descriptively and reported with seed coverage where relevant. An auxiliary
performance-control meta-model diagnostic is reported in Appendix~\ref{app:meta-model-diagnostic}.

\paragraph{Interpretable rule extraction.}
To characterize high-multiplicity regions in terms of original features, shallow decision
trees with maximum depth 3 are used to extract human-readable rules. Trees are fitted on
the original feature space, and positive leaf nodes are converted into conjunctive rules.
The primary target label is HH membership.

Two variants are used. First, global HH rules are fitted to predict membership in the full
HH set of a run. These rules describe broader subgroups enriched for HH observations but
may mix several distinct hotspot regions. Second, component-level rules are fitted for
connected HH components. In this variant, membership in a retained component is the
positive class and all other test observations are the negative class. Component-level rules
therefore aim to describe compact high-multiplicity regions rather than the full HH set at
once.

Rules are evaluated using support, purity, recall, and lift. For global HH rules, purity is
$P(\mathrm{HH}\mid\mathrm{rule\ region})$, and lift is measured relative to the HH base
rate. For component-level rules, purity is
$P(\mathrm{component}\mid\mathrm{rule\ region})$, and lift is measured relative to the
base rate of the corresponding component. The main rule analysis is reported for COMPAS,
where the hotspot structure is clearest and retained components are sufficiently large for
meaningful rule extraction.

\subsection{Chapter Summary}
\label{sec:framework-summary}

\label{sec:meth-assumptions}
The framework relies on four main assumptions. First, the preprocessed feature space must
induce a meaningful notion of local similarity. Second, the finite candidate pool and fixed
top-$K$ selection are treated as a practical approximation of the Rashomon set. Third, the
analysis is descriptive rather than causal: hotspots identify where predictive multiplicity
concentrates, but do not by themselves explain the data-generating mechanisms behind
that disagreement. Fourth, the framework is developed for binary probabilistic
classification and would require adaptation for regression or multiclass settings.

\begin{table}[htbp]
\centering
\small
\caption{Overview of the main analysis goals, methods, and corresponding result sections.}
\label{tab:framework-method-result-map}
\begin{tabular}{p{0.30\textwidth}p{0.34\textwidth}p{0.25\textwidth}}
\toprule
\textbf{Analysis goal} & \textbf{Main method} & \textbf{Result section} \\
\midrule
Measure where near-optimal models disagree & Predictive variance across $\mathcal{R}_K$ & Section~\ref{sec:results-metrics} \\
Test whether disagreement is spatially clustered & Feature-space $k$NN graph, Moran's $I$, prediction-permutation null & Sections~\ref{sec:results-metrics}--\ref{sec:results-null} \\
Identify local high-multiplicity regions & LISA HH labels and connected HH components & Section~\ref{sec:results-spatial} \\
Assess stability and robustness & Fixed-test-set stability, $K$ and $k$ sensitivity, calibration, alternative graphs & Sections~\ref{sec:results-spatial} and~\ref{sec:results-robustness} \\
Summarize associations with modeling choices & Global vs per-family sets, family importance, within-family $V_m$ decomposition & Sections~\ref{sec:results-global-vs-family} and~\ref{sec:results-hp} \\
Validate and interpret hotspots & Synthetic ground-truth designs and shallow-tree rule extraction & Sections~\ref{sec:results-synthetic}--\ref{sec:results-rules} \\
\bottomrule
\end{tabular}
\end{table}
